%%%%%%%%%%%%%%%%%%%%%%%%%%%%%%%%%%%%%%%%%%%%%%%%%%%%%%%%%%%%
% 12.6 RECURSION THEORY
% The Composition of Advanced Mathematics
%%%%%%%%%%%%%%%%%%%%%%%%%%%%%%%%%%%%%%%%%%%%%%%%%%%%%%%%%%%%

\section{Recursion Theory}
\label{sec:recursion-theory}

\prerequisite{
Computability theory, Turing machines, computable and partial computable
functions, computably enumerable sets, reductions, oracle computation,
formal arithmetic, mathematical induction, and elementary number theory.
}

\learningobjectives{
By the end of this section, the reader should be able to:
\begin{itemize}
    \item explain the relationship between recursion theory and
          computability theory;
    \item understand primitive recursive functions;
    \item recognize the basic initial functions;
    \item construct functions using composition;
    \item construct functions using primitive recursion;
    \item prove basic arithmetic functions primitive recursive;
    \item recognize limitations of primitive recursion;
    \item understand the Ackermann function conceptually;
    \item distinguish primitive recursive and general computable functions;
    \item understand minimization;
    \item define partial recursive functions;
    \item recognize total recursive functions;
    \item understand effective numberings of computable functions;
    \item recognize indices of programs;
    \item understand universal partial computable functions;
    \item recognize acceptable numberings conceptually;
    \item understand the \(s\)-\(m\)-\(n\) theorem;
    \item recognize parameterization of programs;
    \item understand Kleene's normal form theorem conceptually;
    \item recognize computation predicates;
    \item understand fixed-point phenomena;
    \item state Kleene's recursion theorem;
    \item recognize self-reproducing and self-referential programs;
    \item understand computably enumerable sets structurally;
    \item recognize enumeration operators conceptually;
    \item distinguish computable, c.e., and co-c.e.\ sets;
    \item understand index sets;
    \item recognize productive sets;
    \item recognize creative sets;
    \item understand many-one completeness for c.e.\ sets;
    \item recognize simple sets conceptually;
    \item understand immune sets conceptually;
    \item recognize computable inseparability;
    \item distinguish many-one and Turing reducibility;
    \item understand one-one reducibility conceptually;
    \item recognize truth-table reducibility conceptually;
    \item understand Turing equivalence;
    \item define Turing degrees;
    \item recognize the computable degree;
    \item understand the halting degree;
    \item understand the Turing jump;
    \item recognize iterated jumps;
    \item understand relative computability;
    \item understand relativization of computability theorems;
    \item recognize Post's problem;
    \item understand intermediate c.e.\ degrees conceptually;
    \item recognize priority arguments conceptually;
    \item understand requirements and injury;
    \item recognize finite-injury constructions;
    \item understand the arithmetic hierarchy;
    \item distinguish \(\Sigma^0_n\), \(\Pi^0_n\), and \(\Delta^0_n\);
    \item understand bounded versus unbounded quantifiers;
    \item recognize Post's theorem;
    \item connect jumps with the arithmetic hierarchy;
    \item understand complete sets at hierarchy levels;
    \item recognize computable approximations;
    \item understand limit computability conceptually;
    \item recognize the limit lemma;
    \item understand enumeration degrees conceptually;
    \item recognize degree structures as partial orders;
    \item understand why there is no single linear scale of
          noncomputability;
    \item connect recursion theory with computability, proof theory,
          arithmetic, logic, type theory, and foundations.
\end{itemize}
}

%%%%%%%%%%%%%%%%%%%%%%%%%%%%%%%%%%%%%%%%%%%%%%%%%%%%%%%%%%%%
\subsection{What Is Recursion Theory?}
\label{subsec:what-is-recursion-theory}
%%%%%%%%%%%%%%%%%%%%%%%%%%%%%%%%%%%%%%%%%%%%%%%%%%%%%%%%%%%%

Recursion theory is the structural study of computability.

Computability theory asks whether a function or set is computable.

Recursion theory goes further and studies:

\[
\boxed{
\text{how computable functions are generated},
}
\]

\[
\boxed{
\text{how noncomputable sets compare},
}
\]

and

\[
\boxed{
\text{how levels of definability and oracle power are organized}.
}
\]

Historically, the terms \textbf{recursion theory} and
\textbf{computability theory} often refer to overlapping areas.

%%%%%%%%%%%%%%%%%%%%%%%%%%%%%%%%%%%%%%%%%%%%%%%%%%%%%%%%%%%%
\subsection{Functions on the Natural Numbers}
\label{subsec:recursive-functions-natural-numbers}
%%%%%%%%%%%%%%%%%%%%%%%%%%%%%%%%%%%%%%%%%%%%%%%%%%%%%%%%%%%%

Recursion theory often studies functions

\[
\boxed{
f:\mathbb{N}^k\to\mathbb{N}
}
\]

and partial functions

\[
\boxed{
f:\mathbb{N}^k\rightharpoonup\mathbb{N}.
}
\]

Finite strings, programs, formulas, and tuples can be encoded by natural
numbers.

Thus computation over arbitrary finite symbolic data can be represented
arithmetically.

%%%%%%%%%%%%%%%%%%%%%%%%%%%%%%%%%%%%%%%%%%%%%%%%%%%%%%%%%%%%
\subsection{Initial Functions}
\label{subsec:initial-recursive-functions}
%%%%%%%%%%%%%%%%%%%%%%%%%%%%%%%%%%%%%%%%%%%%%%%%%%%%%%%%%%%%

Primitive recursive functions are built from a small collection of initial
functions.

The zero function is

\[
\boxed{
Z(n)=0.
}
\]

The successor function is

\[
\boxed{
S(n)=n+1.
}
\]

Projection functions are

\[
\boxed{
P_i^k(x_1,\ldots,x_k)=x_i,
}
\]

where

\[
1\leq i\leq k.
\]

%%%%%%%%%%%%%%%%%%%%%%%%%%%%%%%%%%%%%%%%%%%%%%%%%%%%%%%%%%%%
\subsection{Composition}
\label{subsec:primitive-recursive-composition}
%%%%%%%%%%%%%%%%%%%%%%%%%%%%%%%%%%%%%%%%%%%%%%%%%%%%%%%%%%%%

Suppose

\[
g:\mathbb{N}^m\to\mathbb{N}
\]

and

\[
h_1,\ldots,h_m:
\mathbb{N}^k\to\mathbb{N}
\]

are known functions.

Their composition is

\[
\boxed{
f(\mathbf{x})
=
g
\left(
h_1(\mathbf{x}),
\ldots,
h_m(\mathbf{x})
\right).
}
\]

Primitive recursive functions are closed under composition.

%%%%%%%%%%%%%%%%%%%%%%%%%%%%%%%%%%%%%%%%%%%%%%%%%%%%%%%%%%%%
\subsection{Primitive Recursion}
\label{subsec:primitive-recursion-definition}
%%%%%%%%%%%%%%%%%%%%%%%%%%%%%%%%%%%%%%%%%%%%%%%%%%%%%%%%%%%%

Primitive recursion defines a function from an initial value and a
successor rule.

Suppose

\[
g:\mathbb{N}^k\to\mathbb{N}
\]

and

\[
h:\mathbb{N}^{k+2}\to\mathbb{N}.
\]

Define \(f\) by

\[
\boxed{
f(\mathbf{x},0)
=
g(\mathbf{x}),
}
\]

and

\[
\boxed{
f(\mathbf{x},n+1)
=
h
\left(
\mathbf{x},
n,
f(\mathbf{x},n)
\right).
}
\]

If \(g\) and \(h\) are primitive recursive, then so is \(f\).

%%%%%%%%%%%%%%%%%%%%%%%%%%%%%%%%%%%%%%%%%%%%%%%%%%%%%%%%%%%%
\subsection{Primitive Recursive Functions}
\label{subsec:primitive-recursive-functions}
%%%%%%%%%%%%%%%%%%%%%%%%%%%%%%%%%%%%%%%%%%%%%%%%%%%%%%%%%%%%

\begin{definitionbox}{Primitive Recursive Function}
The class of primitive recursive functions is the smallest class containing
the zero, successor, and projection functions and closed under composition
and primitive recursion.
\end{definitionbox}

Every primitive recursive function is total.

Thus

\[
\boxed{
\text{primitive recursive}
\Longrightarrow
\text{total computable}.
}
\]

The converse is false.

%%%%%%%%%%%%%%%%%%%%%%%%%%%%%%%%%%%%%%%%%%%%%%%%%%%%%%%%%%%%
\subsection{Addition Is Primitive Recursive}
\label{subsec:addition-primitive-recursive}
%%%%%%%%%%%%%%%%%%%%%%%%%%%%%%%%%%%%%%%%%%%%%%%%%%%%%%%%%%%%

Define addition recursively by

\[
\boxed{
x+0=x,
}
\]

and

\[
\boxed{
x+(n+1)
=
S(x+n).
}
\]

Thus addition is obtained by primitive recursion from projection and
successor functions.

Therefore

\[
\boxed{
\operatorname{add}(x,n)=x+n
}
\]

is primitive recursive.

%%%%%%%%%%%%%%%%%%%%%%%%%%%%%%%%%%%%%%%%%%%%%%%%%%%%%%%%%%%%
\subsection{Worked Example: Primitive Recursive Addition}
\label{subsec:worked-primitive-recursive-addition}
%%%%%%%%%%%%%%%%%%%%%%%%%%%%%%%%%%%%%%%%%%%%%%%%%%%%%%%%%%%%

\begin{workedexamplebox}{Computing \(3+2\)}

Using

\[
x+0=x
\]

and

\[
x+(n+1)=S(x+n),
\]

we obtain

\[
3+1
=
S(3+0)
=
S(3)
=
4.
\]

Then

\[
3+2
=
S(3+1)
=
S(4)
=
5.
\]

\begin{answerbox}
\[
\boxed{
3+2=5.
}
\]
\end{answerbox}

\end{workedexamplebox}

%%%%%%%%%%%%%%%%%%%%%%%%%%%%%%%%%%%%%%%%%%%%%%%%%%%%%%%%%%%%
\subsection{Multiplication Is Primitive Recursive}
\label{subsec:multiplication-primitive-recursive}
%%%%%%%%%%%%%%%%%%%%%%%%%%%%%%%%%%%%%%%%%%%%%%%%%%%%%%%%%%%%

Define

\[
\boxed{
x\cdot0=0,
}
\]

and

\[
\boxed{
x\cdot(n+1)
=
x\cdot n+x.
}
\]

Because addition is already primitive recursive, multiplication is primitive
recursive.

Thus

\[
\boxed{
\operatorname{mult}(x,n)=xn
}
\]

belongs to the primitive recursive class.

%%%%%%%%%%%%%%%%%%%%%%%%%%%%%%%%%%%%%%%%%%%%%%%%%%%%%%%%%%%%
\subsection{Exponentiation Is Primitive Recursive}
\label{subsec:exponentiation-primitive-recursive}
%%%%%%%%%%%%%%%%%%%%%%%%%%%%%%%%%%%%%%%%%%%%%%%%%%%%%%%%%%%%

Define

\[
\boxed{
x^0=1,
}
\]

and

\[
\boxed{
x^{n+1}
=
x^n x.
}
\]

Since multiplication is primitive recursive,

\[
\boxed{
\operatorname{exp}(x,n)=x^n
}
\]

is also primitive recursive.

Thus the class already contains rapidly growing arithmetic functions.

%%%%%%%%%%%%%%%%%%%%%%%%%%%%%%%%%%%%%%%%%%%%%%%%%%%%%%%%%%%%
\subsection{Factorial Is Primitive Recursive}
\label{subsec:factorial-primitive-recursive}
%%%%%%%%%%%%%%%%%%%%%%%%%%%%%%%%%%%%%%%%%%%%%%%%%%%%%%%%%%%%

Factorial is defined by

\[
\boxed{
0!=1,
}
\]

and

\[
\boxed{
(n+1)!
=
(n+1)n!.
}
\]

Therefore factorial is primitive recursive.

%%%%%%%%%%%%%%%%%%%%%%%%%%%%%%%%%%%%%%%%%%%%%%%%%%%%%%%%%%%%
\subsection{Truncated Subtraction}
\label{subsec:truncated-subtraction}
%%%%%%%%%%%%%%%%%%%%%%%%%%%%%%%%%%%%%%%%%%%%%%%%%%%%%%%%%%%%

Define predecessor by

\[
\boxed{
\operatorname{pred}(0)=0,
}
\]

and

\[
\boxed{
\operatorname{pred}(n+1)=n.
}
\]

Then define truncated subtraction

\[
\boxed{
x\dotminus0=x,
}
\]

\[
\boxed{
x\dotminus(n+1)
=
\operatorname{pred}(x\dotminus n).
}
\]

Thus

\[
\boxed{
x\dotminus y
=
\max(x-y,0)
}
\]

is primitive recursive.

%%%%%%%%%%%%%%%%%%%%%%%%%%%%%%%%%%%%%%%%%%%%%%%%%%%%%%%%%%%%
\subsection{Primitive Recursive Predicates}
\label{subsec:primitive-recursive-predicates}
%%%%%%%%%%%%%%%%%%%%%%%%%%%%%%%%%%%%%%%%%%%%%%%%%%%%%%%%%%%%

A predicate

\[
R(\mathbf{x})
\]

is primitive recursive if its characteristic function is primitive
recursive.

For example,

\[
\boxed{
x\leq y
}
\]

can be represented using truncated subtraction.

One may define

\[
x\leq y
\iff
x\dotminus y=0.
\]

Hence order relations are primitive recursive.

%%%%%%%%%%%%%%%%%%%%%%%%%%%%%%%%%%%%%%%%%%%%%%%%%%%%%%%%%%%%
\subsection{Bounded Search}
\label{subsec:bounded-search-recursion}
%%%%%%%%%%%%%%%%%%%%%%%%%%%%%%%%%%%%%%%%%%%%%%%%%%%%%%%%%%%%

Primitive recursion can implement bounded search.

Suppose

\[
R(x,y)
\]

is primitive recursive.

One may search

\[
y=0,1,\ldots,n
\]

for the least value satisfying \(R(x,y)\).

Because the search space is finite, the resulting bounded-search function
remains primitive recursive.

%%%%%%%%%%%%%%%%%%%%%%%%%%%%%%%%%%%%%%%%%%%%%%%%%%%%%%%%%%%%
\subsection{Bounded Quantification}
\label{subsec:bounded-quantification-recursion}
%%%%%%%%%%%%%%%%%%%%%%%%%%%%%%%%%%%%%%%%%%%%%%%%%%%%%%%%%%%%

If predicate \(R(x,y)\) is primitive recursive, then bounded statements such
as

\[
\boxed{
\exists y\leq n\,R(x,y)
}
\]

and

\[
\boxed{
\forall y\leq n\,R(x,y)
}
\]

are primitive recursive predicates.

Finite search can be encoded using recursion.

%%%%%%%%%%%%%%%%%%%%%%%%%%%%%%%%%%%%%%%%%%%%%%%%%%%%%%%%%%%%
\subsection{All Primitive Recursive Functions Halt}
\label{subsec:primitive-recursive-totality}
%%%%%%%%%%%%%%%%%%%%%%%%%%%%%%%%%%%%%%%%%%%%%%%%%%%%%%%%%%%%

Every primitive recursive function is defined using explicitly bounded
recursive descent through natural numbers.

Consequently,

\[
\boxed{
\text{every primitive recursive computation terminates}.
}
\]

This makes the class mathematically convenient.

But it also makes the class too small to capture all total computable
functions.

%%%%%%%%%%%%%%%%%%%%%%%%%%%%%%%%%%%%%%%%%%%%%%%%%%%%%%%%%%%%
\subsection{Ackermann-Type Growth}
\label{subsec:ackermann-function}
%%%%%%%%%%%%%%%%%%%%%%%%%%%%%%%%%%%%%%%%%%%%%%%%%%%%%%%%%%%%

The Ackermann function is a classical example of a total computable function
that is not primitive recursive.

One common form is defined recursively by

\[
A(0,n)
=
n+1,
\]

\[
A(m+1,0)
=
A(m,1),
\]

and

\[
A(m+1,n+1)
=
A
\left(
m,
A(m+1,n)
\right).
\]

Its growth eventually dominates every primitive recursive function in an
appropriate sense.

%%%%%%%%%%%%%%%%%%%%%%%%%%%%%%%%%%%%%%%%%%%%%%%%%%%%%%%%%%%%
\subsection{Primitive Recursive Is Not All Computable}
\label{subsec:primitive-recursive-not-all}
%%%%%%%%%%%%%%%%%%%%%%%%%%%%%%%%%%%%%%%%%%%%%%%%%%%%%%%%%%%%

The Ackermann example establishes

\[
\boxed{
\text{primitive recursive}
\subsetneq
\text{total computable}.
}
\]

Thus an additional operation is required to capture general computability.

One standard extension is unbounded minimization.

%%%%%%%%%%%%%%%%%%%%%%%%%%%%%%%%%%%%%%%%%%%%%%%%%%%%%%%%%%%%
\subsection{Minimization Operator}
\label{subsec:minimization-operator}
%%%%%%%%%%%%%%%%%%%%%%%%%%%%%%%%%%%%%%%%%%%%%%%%%%%%%%%%%%%%

Suppose

\[
g(\mathbf{x},y)
\]

is a computable function.

Define

\[
\boxed{
\mu y
\,
[g(\mathbf{x},y)=0]
}
\]

to be the least \(y\) such that

\[
g(\mathbf{x},y)=0.
\]

The search begins

\[
0,1,2,\ldots
\]

and stops at the first successful value.

If no such \(y\) exists, the search never terminates.

%%%%%%%%%%%%%%%%%%%%%%%%%%%%%%%%%%%%%%%%%%%%%%%%%%%%%%%%%%%%
\subsection{Partiality From Minimization}
\label{subsec:minimization-partiality}
%%%%%%%%%%%%%%%%%%%%%%%%%%%%%%%%%%%%%%%%%%%%%%%%%%%%%%%%%%%%

Unlike bounded search, minimization is unbounded.

Therefore

\[
\boxed{
\mu y
\,
[g(\mathbf{x},y)=0]
}
\]

may be undefined.

This creates partial computable functions naturally.

Thus nontermination enters the recursive-function framework directly.

%%%%%%%%%%%%%%%%%%%%%%%%%%%%%%%%%%%%%%%%%%%%%%%%%%%%%%%%%%%%
\subsection{Partial Recursive Functions}
\label{subsec:partial-recursive-functions}
%%%%%%%%%%%%%%%%%%%%%%%%%%%%%%%%%%%%%%%%%%%%%%%%%%%%%%%%%%%%

\begin{definitionbox}{Partial Recursive Function}
The partial recursive functions are generated from the initial functions
using composition, primitive recursion, and unbounded minimization.
\end{definitionbox}

The resulting class coincides with the class of partial Turing-computable
functions.

Thus

\[
\boxed{
\text{partial recursive}
=
\text{partial computable}.
}
\]

%%%%%%%%%%%%%%%%%%%%%%%%%%%%%%%%%%%%%%%%%%%%%%%%%%%%%%%%%%%%
\subsection{Total Recursive Functions}
\label{subsec:total-recursive-functions}
%%%%%%%%%%%%%%%%%%%%%%%%%%%%%%%%%%%%%%%%%%%%%%%%%%%%%%%%%%%%

A partial recursive function that is defined on every input is called total
recursive.

Thus

\[
\boxed{
\text{recursive function}
}
\]

in classical terminology often means

\[
\boxed{
\text{total computable function}.
}
\]

Terminology varies, so the intended meaning should always be stated.

%%%%%%%%%%%%%%%%%%%%%%%%%%%%%%%%%%%%%%%%%%%%%%%%%%%%%%%%%%%%
\subsection{Partial Versus Total Recursive Functions}
\label{subsec:partial-total-recursive}
%%%%%%%%%%%%%%%%%%%%%%%%%%%%%%%%%%%%%%%%%%%%%%%%%%%%%%%%%%%%

The difference is domain.

A total recursive function satisfies

\[
\boxed{
\operatorname{dom}(f)
=
\mathbb{N}^k.
}
\]

A partial recursive function may satisfy

\[
\boxed{
\operatorname{dom}(f)
\subsetneq
\mathbb{N}^k.
}
\]

Undefinedness corresponds computationally to nontermination.

%%%%%%%%%%%%%%%%%%%%%%%%%%%%%%%%%%%%%%%%%%%%%%%%%%%%%%%%%%%%
\subsection{Numbering Computable Functions}
\label{subsec:numbering-computable-functions}
%%%%%%%%%%%%%%%%%%%%%%%%%%%%%%%%%%%%%%%%%%%%%%%%%%%%%%%%%%%%

Because programs have finite descriptions, partial computable functions can
be effectively enumerated:

\[
\boxed{
\varphi_0,
\varphi_1,
\varphi_2,
\ldots.
}
\]

The number \(e\) is called an index of the program or partial computable
function \(\varphi_e\).

A function may have many different indices.

%%%%%%%%%%%%%%%%%%%%%%%%%%%%%%%%%%%%%%%%%%%%%%%%%%%%%%%%%%%%
\subsection{Indices Are Not Unique}
\label{subsec:indices-not-unique}
%%%%%%%%%%%%%%%%%%%%%%%%%%%%%%%%%%%%%%%%%%%%%%%%%%%%%%%%%%%%

Two distinct programs can compute the same function.

Thus it is possible that

\[
\boxed{
e\neq d
}
\]

while

\[
\boxed{
\varphi_e
=
\varphi_d.
}
\]

An index describes a program representation, not a unique mathematical
function.

%%%%%%%%%%%%%%%%%%%%%%%%%%%%%%%%%%%%%%%%%%%%%%%%%%%%%%%%%%%%
\subsection{Universal Partial Computable Function}
\label{subsec:universal-partial-computable-function}
%%%%%%%%%%%%%%%%%%%%%%%%%%%%%%%%%%%%%%%%%%%%%%%%%%%%%%%%%%%%

There exists a partial computable function

\[
\boxed{
U(e,x)
}
\]

such that

\[
\boxed{
U(e,x)
=
\varphi_e(x)
}
\]

whenever \(\varphi_e(x)\) is defined.

Thus one computable mechanism can simulate every indexed partial computable
function.

%%%%%%%%%%%%%%%%%%%%%%%%%%%%%%%%%%%%%%%%%%%%%%%%%%%%%%%%%%%%
\subsection{The \(s\)-\(m\)-\(n\) Theorem}
\label{subsec:s-m-n-theorem}
%%%%%%%%%%%%%%%%%%%%%%%%%%%%%%%%%%%%%%%%%%%%%%%%%%%%%%%%%%%%

The \(s\)-\(m\)-\(n\) theorem formalizes effective parameter substitution
into programs.

Suppose

\[
\varphi_e(x,y)
\]

is a partial computable function of two arguments.

There exists a total computable function \(s\) such that

\[
\boxed{
\varphi_{s(e,x)}(y)
=
\varphi_e(x,y).
}
\]

Thus one can effectively produce a program with parameter \(x\) already
built into its code.

%%%%%%%%%%%%%%%%%%%%%%%%%%%%%%%%%%%%%%%%%%%%%%%%%%%%%%%%%%%%
\subsection{Meaning of Parameterization}
\label{subsec:meaning-s-m-n}
%%%%%%%%%%%%%%%%%%%%%%%%%%%%%%%%%%%%%%%%%%%%%%%%%%%%%%%%%%%%

Suppose program \(P\) takes inputs

\[
(x,y).
\]

The \(s\)-\(m\)-\(n\) theorem says there is an effective transformation that
takes:

\[
\boxed{
\text{code of }P
}
\]

and

\[
\boxed{
\text{fixed value }x
}
\]

and produces code for a new program computing

\[
\boxed{
y\mapsto P(x,y).
}
\]

This is a formal version of partial evaluation.

%%%%%%%%%%%%%%%%%%%%%%%%%%%%%%%%%%%%%%%%%%%%%%%%%%%%%%%%%%%%
\subsection{Worked Example: Hard-Coding a Parameter}
\label{subsec:worked-hard-coded-parameter}
%%%%%%%%%%%%%%%%%%%%%%%%%%%%%%%%%%%%%%%%%%%%%%%%%%%%%%%%%%%%

\begin{workedexamplebox}{Specializing Addition}

Suppose a program computes

\[
f(x,y)=x+y.
\]

Fix

\[
x=5.
\]

Parameterization produces a new program computing

\[
g(y)=5+y.
\]

Thus the general two-input program has been specialized to one fixed first
argument.

\begin{answerbox}
\[
\boxed{
g(y)=f(5,y)=5+y.
}
\]
\end{answerbox}

The \(s\)-\(m\)-\(n\) theorem states that this specialization can be
performed effectively at the level of program indices.

\end{workedexamplebox}

%%%%%%%%%%%%%%%%%%%%%%%%%%%%%%%%%%%%%%%%%%%%%%%%%%%%%%%%%%%%
\subsection{Kleene Normal Form}
\label{subsec:kleene-normal-form}
%%%%%%%%%%%%%%%%%%%%%%%%%%%%%%%%%%%%%%%%%%%%%%%%%%%%%%%%%%%%

Kleene's normal form theorem gives a uniform arithmetic representation of
partial computable functions.

Conceptually, there exist primitive recursive objects that encode:

\[
\boxed{
\text{a valid computation history},
}
\]

and

\[
\boxed{
\text{the output of that computation}.
}
\]

A partial computable function can then be represented using minimization
over valid computation codes.

%%%%%%%%%%%%%%%%%%%%%%%%%%%%%%%%%%%%%%%%%%%%%%%%%%%%%%%%%%%%
\subsection{Computation Predicate}
\label{subsec:kleene-computation-predicate}
%%%%%%%%%%%%%%%%%%%%%%%%%%%%%%%%%%%%%%%%%%%%%%%%%%%%%%%%%%%%

A primitive recursive relation

\[
T(e,x,s)
\]

can be interpreted conceptually as:

\[
\boxed{
s
\text{ codes a halting computation of program }e
\text{ on input }x.
}
\]

Then one searches for the least valid computation history.

This provides an arithmetic representation of computation.

%%%%%%%%%%%%%%%%%%%%%%%%%%%%%%%%%%%%%%%%%%%%%%%%%%%%%%%%%%%%
\subsection{Normal Form Schema}
\label{subsec:kleene-normal-form-schema}
%%%%%%%%%%%%%%%%%%%%%%%%%%%%%%%%%%%%%%%%%%%%%%%%%%%%%%%%%%%%

Schematically, one may express a partial computable function as

\[
\boxed{
\varphi_e(x)
=
U
\left(
\mu s\,
T(e,x,s)
\right),
}
\]

where the symbols \(T\) and \(U\) here represent suitable primitive
recursive coding functions.

The exact standard notation varies by presentation.

The important point is:

\[
\boxed{
\text{general computation}
=
\text{primitive recursive verification}
+
\text{unbounded search}.
}
\]

%%%%%%%%%%%%%%%%%%%%%%%%%%%%%%%%%%%%%%%%%%%%%%%%%%%%%%%%%%%%
\subsection{Kleene's Recursion Theorem}
\label{subsec:kleene-recursion-theorem}
%%%%%%%%%%%%%%%%%%%%%%%%%%%%%%%%%%%%%%%%%%%%%%%%%%%%%%%%%%%%

The recursion theorem establishes fixed points for effective program
transformations.

One standard form states:

if

\[
f:\mathbb{N}\to\mathbb{N}
\]

is total computable, then there exists an index \(e\) such that

\[
\boxed{
\varphi_e
=
\varphi_{f(e)}.
}
\]

Thus a program can behave exactly like a computably transformed version of
its own description.

%%%%%%%%%%%%%%%%%%%%%%%%%%%%%%%%%%%%%%%%%%%%%%%%%%%%%%%%%%%%
\subsection{Why the Recursion Theorem Is Surprising}
\label{subsec:why-recursion-theorem}
%%%%%%%%%%%%%%%%%%%%%%%%%%%%%%%%%%%%%%%%%%%%%%%%%%%%%%%%%%%%

The theorem permits controlled self-reference without giving a program an
external copy of its own source code.

A program may effectively behave according to a transformation involving
its own index.

This supports formal constructions of:

\[
\boxed{
\text{self-reproducing programs},
}
\]

\[
\boxed{
\text{self-referential machines},
}
\]

and

\[
\boxed{
\text{fixed points of program transformations}.
}
\]

%%%%%%%%%%%%%%%%%%%%%%%%%%%%%%%%%%%%%%%%%%%%%%%%%%%%%%%%%%%%
\subsection{Self-Reproducing Programs}
\label{subsec:self-reproducing-programs}
%%%%%%%%%%%%%%%%%%%%%%%%%%%%%%%%%%%%%%%%%%%%%%%%%%%%%%%%%%%%

A self-reproducing program outputs a description of itself.

The recursion theorem explains why such behavior is possible in any
sufficiently expressive effective programming system.

The result is mathematical and does not depend on one specific programming
language.

%%%%%%%%%%%%%%%%%%%%%%%%%%%%%%%%%%%%%%%%%%%%%%%%%%%%%%%%%%%%
\subsection{Recursion Theorem and Gödel Self-Reference}
\label{subsec:recursion-theorem-godel}
%%%%%%%%%%%%%%%%%%%%%%%%%%%%%%%%%%%%%%%%%%%%%%%%%%%%%%%%%%%%

The recursion theorem and Gödel's diagonal lemma are closely related
fixed-point phenomena.

Both allow an object to refer indirectly to its own description.

In proof theory:

\[
\boxed{
\text{sentences refer to their own codes}.
}
\]

In recursion theory:

\[
\boxed{
\text{programs act using their own indices}.
}
\]

%%%%%%%%%%%%%%%%%%%%%%%%%%%%%%%%%%%%%%%%%%%%%%%%%%%%%%%%%%%%
\subsection{Computably Enumerable Sets}
\label{subsec:recursion-theory-ce-sets}
%%%%%%%%%%%%%%%%%%%%%%%%%%%%%%%%%%%%%%%%%%%%%%%%%%%%%%%%%%%%

A set

\[
A\subseteq\mathbb{N}
\]

is computably enumerable if some partial computable function recognizes its
members.

Equivalent descriptions include:

\[
\boxed{
A
=
\operatorname{dom}(\varphi_e)
}
\]

for some \(e\), or

\[
\boxed{
A
=
\operatorname{range}(f)
}
\]

for a suitable computable enumeration, with standard qualifications for the
empty set.

%%%%%%%%%%%%%%%%%%%%%%%%%%%%%%%%%%%%%%%%%%%%%%%%%%%%%%%%%%%%
\subsection{Domain Representation}
\label{subsec:ce-domain-representation}
%%%%%%%%%%%%%%%%%%%%%%%%%%%%%%%%%%%%%%%%%%%%%%%%%%%%%%%%%%%%

For every c.e.\ set \(A\), there exists an index \(e\) such that

\[
\boxed{
A
=
W_e
=
\operatorname{dom}(\varphi_e).
}
\]

The notation

\[
W_e
\]

is standard in recursion theory for the \(e\)-th computably enumerable set.

%%%%%%%%%%%%%%%%%%%%%%%%%%%%%%%%%%%%%%%%%%%%%%%%%%%%%%%%%%%%
\subsection{Stage Approximations}
\label{subsec:ce-stage-approximations}
%%%%%%%%%%%%%%%%%%%%%%%%%%%%%%%%%%%%%%%%%%%%%%%%%%%%%%%%%%%%

A c.e.\ set may be enumerated gradually.

Write

\[
\boxed{
W_{e,0}
\subseteq
W_{e,1}
\subseteq
W_{e,2}
\subseteq
\cdots
}
\]

for finite-stage approximations.

Then

\[
\boxed{
W_e
=
\bigcup_s
W_{e,s}.
}
\]

Membership may appear at a finite stage, but nonmembership may never be
confirmed.

%%%%%%%%%%%%%%%%%%%%%%%%%%%%%%%%%%%%%%%%%%%%%%%%%%%%%%%%%%%%
\subsection{Computable Sets as Double Enumeration}
\label{subsec:computable-double-enumeration}
%%%%%%%%%%%%%%%%%%%%%%%%%%%%%%%%%%%%%%%%%%%%%%%%%%%%%%%%%%%%

A set \(A\) is computable if and only if both

\[
\boxed{
A
}
\]

and

\[
\boxed{
\overline{A}
}
\]

are computably enumerable.

By dovetailing their enumerations, every input eventually appears on exactly
one side.

Thus membership becomes decidable.

%%%%%%%%%%%%%%%%%%%%%%%%%%%%%%%%%%%%%%%%%%%%%%%%%%%%%%%%%%%%
\subsection{Co-c.e.\ Sets}
\label{subsec:co-ce-sets}
%%%%%%%%%%%%%%%%%%%%%%%%%%%%%%%%%%%%%%%%%%%%%%%%%%%%%%%%%%%%

A set \(A\) is co-c.e.\ if

\[
\boxed{
\overline{A}
}
\]

is computably enumerable.

Hence:

\[
\boxed{
A
\text{ computable}
\iff
A
\text{ is c.e.\ and co-c.e.}
}
\]

This is one of the basic structural facts of recursion theory.

%%%%%%%%%%%%%%%%%%%%%%%%%%%%%%%%%%%%%%%%%%%%%%%%%%%%%%%%%%%%
\subsection{The Halting Set}
\label{subsec:recursion-halting-set}
%%%%%%%%%%%%%%%%%%%%%%%%%%%%%%%%%%%%%%%%%%%%%%%%%%%%%%%%%%%%

Define the diagonal halting set

\[
\boxed{
K
=
\{
e:
\varphi_e(e)
\downarrow
\},
}
\]

where

\[
\downarrow
\]

means that the computation halts.

The set \(K\) is computably enumerable but not computable.

It is a central complete c.e.\ set.

%%%%%%%%%%%%%%%%%%%%%%%%%%%%%%%%%%%%%%%%%%%%%%%%%%%%%%%%%%%%
\subsection{Many-One Completeness}
\label{subsec:ce-many-one-completeness}
%%%%%%%%%%%%%%%%%%%%%%%%%%%%%%%%%%%%%%%%%%%%%%%%%%%%%%%%%%%%

A c.e.\ set \(A\) is many-one complete for the c.e.\ sets if:

\[
A
\text{ is c.e.}
\]

and for every c.e.\ set \(B\),

\[
\boxed{
B
\leq_m
A.
}
\]

Such a set is maximally difficult among c.e.\ sets under many-one
reducibility.

The halting set provides a standard example.

%%%%%%%%%%%%%%%%%%%%%%%%%%%%%%%%%%%%%%%%%%%%%%%%%%%%%%%%%%%%
\subsection{Index Sets}
\label{subsec:index-sets}
%%%%%%%%%%%%%%%%%%%%%%%%%%%%%%%%%%%%%%%%%%%%%%%%%%%%%%%%%%%%

An index set is a set of program indices determined only by the function or
set computed.

If

\[
\varphi_e
=
\varphi_d,
\]

then an index set \(I\) satisfies

\[
\boxed{
e\in I
\iff
d\in I.
}
\]

Thus membership depends on semantic behavior rather than syntactic code.

Rice's theorem concerns nontrivial computable properties of such indices.

%%%%%%%%%%%%%%%%%%%%%%%%%%%%%%%%%%%%%%%%%%%%%%%%%%%%%%%%%%%%
\subsection{Productive Sets}
\label{subsec:productive-sets}
%%%%%%%%%%%%%%%%%%%%%%%%%%%%%%%%%%%%%%%%%%%%%%%%%%%%%%%%%%%%

A set \(A\) is productive, roughly, if there exists a computable procedure
that defeats every attempted c.e.\ enumeration contained within \(A\).

More precisely, one seeks a partial computable function \(p\) such that if

\[
W_e\subseteq A,
\]

then

\[
\boxed{
p(e)\in A\setminus W_e.
}
\]

Thus no c.e.\ subset can exhaust \(A\).

%%%%%%%%%%%%%%%%%%%%%%%%%%%%%%%%%%%%%%%%%%%%%%%%%%%%%%%%%%%%
\subsection{Creative Sets}
\label{subsec:creative-sets}
%%%%%%%%%%%%%%%%%%%%%%%%%%%%%%%%%%%%%%%%%%%%%%%%%%%%%%%%%%%%

A c.e.\ set \(C\) is creative if its complement

\[
\boxed{
\overline{C}
}
\]

is productive.

Creative sets represent strongly noncomputable c.e.\ sets.

Classically, creative sets are closely connected with many-one complete
c.e.\ sets.

%%%%%%%%%%%%%%%%%%%%%%%%%%%%%%%%%%%%%%%%%%%%%%%%%%%%%%%%%%%%
\subsection{Computable Inseparability}
\label{subsec:computable-inseparability}
%%%%%%%%%%%%%%%%%%%%%%%%%%%%%%%%%%%%%%%%%%%%%%%%%%%%%%%%%%%%

Two disjoint sets

\[
A,B\subseteq\mathbb{N}
\]

are computably inseparable if there is no computable set \(C\) such that

\[
\boxed{
A\subseteq C
}
\]

and

\[
\boxed{
B\cap C=\varnothing.
}
\]

Thus no decidable set separates all elements of \(A\) from all elements of
\(B\).

%%%%%%%%%%%%%%%%%%%%%%%%%%%%%%%%%%%%%%%%%%%%%%%%%%%%%%%%%%%%
\subsection{Why Inseparability Matters}
\label{subsec:why-inseparability-matters}
%%%%%%%%%%%%%%%%%%%%%%%%%%%%%%%%%%%%%%%%%%%%%%%%%%%%%%%%%%%%

Computable inseparability provides a refined form of undecidability.

Instead of one set being noncomputable, two disjoint c.e.\ sets may be so
intertwined that no computable boundary separates them.

Such ideas occur in logic and incompleteness constructions.

%%%%%%%%%%%%%%%%%%%%%%%%%%%%%%%%%%%%%%%%%%%%%%%%%%%%%%%%%%%%
\subsection{Immune Sets}
\label{subsec:immune-sets}
%%%%%%%%%%%%%%%%%%%%%%%%%%%%%%%%%%%%%%%%%%%%%%%%%%%%%%%%%%%%

An infinite set \(A\) is immune if it contains no infinite c.e.\ subset.

Thus

\[
\boxed{
A
\text{ infinite}
}
\]

but every c.e.\ set

\[
W_e\subseteq A
\]

is finite.

Immune sets illustrate how an infinite set can resist effective
enumeration.

%%%%%%%%%%%%%%%%%%%%%%%%%%%%%%%%%%%%%%%%%%%%%%%%%%%%%%%%%%%%
\subsection{Simple Sets}
\label{subsec:simple-sets}
%%%%%%%%%%%%%%%%%%%%%%%%%%%%%%%%%%%%%%%%%%%%%%%%%%%%%%%%%%%%

A c.e.\ set \(A\) is simple if:

\[
\boxed{
\overline{A}
\text{ is infinite}
}
\]

and

\[
\boxed{
\overline{A}
\text{ contains no infinite c.e.\ subset}.
}
\]

Thus the complement is immune.

Simple sets show that not every noncomputable c.e.\ set is complete.

%%%%%%%%%%%%%%%%%%%%%%%%%%%%%%%%%%%%%%%%%%%%%%%%%%%%%%%%%%%%
\subsection{Why Simple Sets Are Important}
\label{subsec:why-simple-sets-important}
%%%%%%%%%%%%%%%%%%%%%%%%%%%%%%%%%%%%%%%%%%%%%%%%%%%%%%%%%%%%

The existence of simple sets demonstrates a hierarchy inside the c.e.\
sets.

There are c.e.\ sets that are:

\[
\boxed{
\text{noncomputable}
}
\]

but not maximally difficult under standard completeness notions.

This motivated the study of intermediate degrees of unsolvability.

%%%%%%%%%%%%%%%%%%%%%%%%%%%%%%%%%%%%%%%%%%%%%%%%%%%%%%%%%%%%
\subsection{Reducibility as Comparison of Information}
\label{subsec:reducibility-information}
%%%%%%%%%%%%%%%%%%%%%%%%%%%%%%%%%%%%%%%%%%%%%%%%%%%%%%%%%%%%

A reducibility relation compares computational information.

If

\[
A\leq B,
\]

then intuitively:

\[
\boxed{
B
\text{ contains enough effective information to solve }A.
}
\]

Different reducibilities impose different restrictions on how that
information may be used.

%%%%%%%%%%%%%%%%%%%%%%%%%%%%%%%%%%%%%%%%%%%%%%%%%%%%%%%%%%%%
\subsection{Many-One Reducibility}
\label{subsec:recursion-many-one-reducibility}
%%%%%%%%%%%%%%%%%%%%%%%%%%%%%%%%%%%%%%%%%%%%%%%%%%%%%%%%%%%%

Recall:

\[
\boxed{
A\leq_m B
}
\]

if there is a total computable function \(f\) such that

\[
\boxed{
x\in A
\iff
f(x)\in B.
}
\]

Only one transformed query to \(B\) is needed.

This is a relatively strong restriction on the reduction.

%%%%%%%%%%%%%%%%%%%%%%%%%%%%%%%%%%%%%%%%%%%%%%%%%%%%%%%%%%%%
\subsection{One-One Reducibility}
\label{subsec:one-one-reducibility}
%%%%%%%%%%%%%%%%%%%%%%%%%%%%%%%%%%%%%%%%%%%%%%%%%%%%%%%%%%%%

A one-one reduction is a mapping reduction in which the reducing function
is injective.

Write

\[
\boxed{
A\leq_1B.
}
\]

Then

\[
\boxed{
A\leq_1B
\Longrightarrow
A\leq_mB.
}
\]

One-one reducibility preserves more structural information than arbitrary
many-one reducibility.

%%%%%%%%%%%%%%%%%%%%%%%%%%%%%%%%%%%%%%%%%%%%%%%%%%%%%%%%%%%%
\subsection{Truth-Table Reducibility}
\label{subsec:truth-table-reducibility}
%%%%%%%%%%%%%%%%%%%%%%%%%%%%%%%%%%%%%%%%%%%%%%%%%%%%%%%%%%%%

Truth-table reducibility permits finitely many oracle questions whose
queries are determined effectively before receiving the oracle answers.

The final answer is then computed from a fixed Boolean combination of the
responses.

Conceptually,

\[
\boxed{
\text{nonadaptive finite oracle access}.
}
\]

%%%%%%%%%%%%%%%%%%%%%%%%%%%%%%%%%%%%%%%%%%%%%%%%%%%%%%%%%%%%
\subsection{Turing Reducibility}
\label{subsec:recursion-turing-reducibility}
%%%%%%%%%%%%%%%%%%%%%%%%%%%%%%%%%%%%%%%%%%%%%%%%%%%%%%%%%%%%

Turing reducibility permits an algorithm to query an oracle repeatedly and
adaptively.

Write

\[
\boxed{
A\leq_TB.
}
\]

The next query may depend on answers to earlier queries.

Therefore Turing reducibility is more flexible than many-one or
truth-table reducibility.

%%%%%%%%%%%%%%%%%%%%%%%%%%%%%%%%%%%%%%%%%%%%%%%%%%%%%%%%%%%%
\subsection{Hierarchy of Reducibilities}
\label{subsec:hierarchy-reducibilities}
%%%%%%%%%%%%%%%%%%%%%%%%%%%%%%%%%%%%%%%%%%%%%%%%%%%%%%%%%%%%

Conceptually,

\[
\boxed{
A\leq_1B
\Longrightarrow
A\leq_mB
}
\]

and standard computability theory gives stronger oracle reducibilities above
these.

In particular,

\[
\boxed{
A\leq_mB
\Longrightarrow
A\leq_TB.
}
\]

The converses need not hold.

%%%%%%%%%%%%%%%%%%%%%%%%%%%%%%%%%%%%%%%%%%%%%%%%%%%%%%%%%%%%
\subsection{Turing Equivalence}
\label{subsec:turing-equivalence}
%%%%%%%%%%%%%%%%%%%%%%%%%%%%%%%%%%%%%%%%%%%%%%%%%%%%%%%%%%%%

Sets \(A\) and \(B\) are Turing equivalent if

\[
\boxed{
A\leq_TB
}
\]

and

\[
\boxed{
B\leq_TA.
}
\]

Write

\[
\boxed{
A\equiv_TB.
}
\]

They possess the same computational power as oracles.

%%%%%%%%%%%%%%%%%%%%%%%%%%%%%%%%%%%%%%%%%%%%%%%%%%%%%%%%%%%%
\subsection{Turing Degrees}
\label{subsec:turing-degrees}
%%%%%%%%%%%%%%%%%%%%%%%%%%%%%%%%%%%%%%%%%%%%%%%%%%%%%%%%%%%%

A Turing degree is an equivalence class under

\[
\equiv_T.
\]

The degree of \(A\) is written

\[
\boxed{
\deg_T(A).
}
\]

Degrees are partially ordered by Turing reducibility.

Thus recursion theory studies a structure of computational information
levels.

%%%%%%%%%%%%%%%%%%%%%%%%%%%%%%%%%%%%%%%%%%%%%%%%%%%%%%%%%%%%
\subsection{The Computable Degree}
\label{subsec:computable-degree}
%%%%%%%%%%%%%%%%%%%%%%%%%%%%%%%%%%%%%%%%%%%%%%%%%%%%%%%%%%%%

Every computable set has the same Turing degree.

This least degree is usually denoted

\[
\boxed{
\mathbf{0}.
}
\]

For every set \(A\),

\[
\boxed{
\mathbf{0}
\leq_T
\deg_T(A).
}
\]

Thus computable information forms the base of the degree structure.

%%%%%%%%%%%%%%%%%%%%%%%%%%%%%%%%%%%%%%%%%%%%%%%%%%%%%%%%%%%%
\subsection{The Halting Degree}
\label{subsec:recursion-halting-degree}
%%%%%%%%%%%%%%%%%%%%%%%%%%%%%%%%%%%%%%%%%%%%%%%%%%%%%%%%%%%%

The degree of the halting set is denoted

\[
\boxed{
\mathbf{0}'.
}
\]

It is strictly above the computable degree:

\[
\boxed{
\mathbf{0}
<
\mathbf{0}'.
}
\]

An oracle for the halting problem can solve strictly more problems than an
ordinary computable procedure.

%%%%%%%%%%%%%%%%%%%%%%%%%%%%%%%%%%%%%%%%%%%%%%%%%%%%%%%%%%%%
\subsection{The Turing Jump}
\label{subsec:recursion-turing-jump}
%%%%%%%%%%%%%%%%%%%%%%%%%%%%%%%%%%%%%%%%%%%%%%%%%%%%%%%%%%%%

Given set \(A\), define its jump \(A'\) as a halting problem for machines
with oracle \(A\).

Conceptually,

\[
\boxed{
A'
=
\{
e:
\varphi_e^A(e)\downarrow
\}.
}
\]

The jump satisfies

\[
\boxed{
A<_TA'.
}
\]

Thus no oracle solves its own associated halting problem.

%%%%%%%%%%%%%%%%%%%%%%%%%%%%%%%%%%%%%%%%%%%%%%%%%%%%%%%%%%%%
\subsection{Iterated Jumps}
\label{subsec:iterated-turing-jumps}
%%%%%%%%%%%%%%%%%%%%%%%%%%%%%%%%%%%%%%%%%%%%%%%%%%%%%%%%%%%%

Starting from the computable degree:

\[
\boxed{
\mathbf{0},
}
\]

one obtains

\[
\boxed{
\mathbf{0}',
}
\]

then

\[
\boxed{
\mathbf{0}'',
}
\]

and so on.

These iterated jumps correspond closely to increasing levels of arithmetic
definability.

%%%%%%%%%%%%%%%%%%%%%%%%%%%%%%%%%%%%%%%%%%%%%%%%%%%%%%%%%%%%
\subsection{Relativization}
\label{subsec:recursion-relativization}
%%%%%%%%%%%%%%%%%%%%%%%%%%%%%%%%%%%%%%%%%%%%%%%%%%%%%%%%%%%%

Many computability definitions and theorems can be relativized to an oracle
\(A\).

Examples include:

\[
\boxed{
A\text{-computable functions},
}
\]

\[
\boxed{
A\text{-c.e.\ sets},
}
\]

and

\[
\boxed{
A\text{-recursive constructions}.
}
\]

The ordinary computability theory is the special case with no noncomputable
oracle.

%%%%%%%%%%%%%%%%%%%%%%%%%%%%%%%%%%%%%%%%%%%%%%%%%%%%%%%%%%%%
\subsection{Relativized Halting Problem}
\label{subsec:relativized-halting-problem}
%%%%%%%%%%%%%%%%%%%%%%%%%%%%%%%%%%%%%%%%%%%%%%%%%%%%%%%%%%%%

Even with oracle \(A\), there is no \(A\)-computable procedure solving the
halting problem for all \(A\)-oracle programs.

That problem is represented by

\[
A'.
\]

Thus undecidability persists at every oracle level.

%%%%%%%%%%%%%%%%%%%%%%%%%%%%%%%%%%%%%%%%%%%%%%%%%%%%%%%%%%%%
\subsection{No Ultimate Halting Oracle}
\label{subsec:no-ultimate-halting-oracle}
%%%%%%%%%%%%%%%%%%%%%%%%%%%%%%%%%%%%%%%%%%%%%%%%%%%%%%%%%%%%

Suppose one gains access to oracle

\[
A'.
\]

This solves the halting problem relative to \(A\).

But then one can form

\[
\boxed{
A''.
}
\]

Therefore there is no final finite stage at which all relative halting
questions disappear.

%%%%%%%%%%%%%%%%%%%%%%%%%%%%%%%%%%%%%%%%%%%%%%%%%%%%%%%%%%%%
\subsection{Post's Problem}
\label{subsec:posts-problem}
%%%%%%%%%%%%%%%%%%%%%%%%%%%%%%%%%%%%%%%%%%%%%%%%%%%%%%%%%%%%

The c.e.\ sets include:

\[
\boxed{
\text{computable sets}
}
\]

of degree \(\mathbf{0}\), and complete c.e.\ sets of degree

\[
\boxed{
\mathbf{0}'.
}
\]

Post asked whether there exists a c.e.\ set \(A\) satisfying

\[
\boxed{
\mathbf{0}
<
\deg_T(A)
<
\mathbf{0}'.
}
\]

That is, is there an intermediate c.e.\ degree?

The answer is yes.

%%%%%%%%%%%%%%%%%%%%%%%%%%%%%%%%%%%%%%%%%%%%%%%%%%%%%%%%%%%%
\subsection{Intermediate c.e.\ Degrees}
\label{subsec:intermediate-ce-degrees}
%%%%%%%%%%%%%%%%%%%%%%%%%%%%%%%%%%%%%%%%%%%%%%%%%%%%%%%%%%%%

The solution of Post's problem showed that there exist c.e.\ sets \(A\)
such that

\[
\boxed{
\mathbf{0}
<
\deg_T(A)
<
\mathbf{0}'.
}
\]

Thus c.e.\ noncomputability is not divided merely into:

\[
\boxed{
\text{computable}
}
\]

and

\[
\boxed{
\text{complete}.
}
\]

There are intermediate computational strengths.

%%%%%%%%%%%%%%%%%%%%%%%%%%%%%%%%%%%%%%%%%%%%%%%%%%%%%%%%%%%%
\subsection{Priority Method}
\label{subsec:priority-method}
%%%%%%%%%%%%%%%%%%%%%%%%%%%%%%%%%%%%%%%%%%%%%%%%%%%%%%%%%%%%

The priority method constructs computably enumerable sets satisfying many
competing requirements simultaneously.

A construction may list requirements

\[
\boxed{
R_0,
R_1,
R_2,\ldots.
}
\]

Higher-priority requirements are protected when conflicts arise.

This method was central to the solution of Post's problem.

%%%%%%%%%%%%%%%%%%%%%%%%%%%%%%%%%%%%%%%%%%%%%%%%%%%%%%%%%%%%
\subsection{Requirements}
\label{subsec:priority-requirements}
%%%%%%%%%%%%%%%%%%%%%%%%%%%%%%%%%%%%%%%%%%%%%%%%%%%%%%%%%%%%

A requirement expresses one property the constructed set must satisfy.

For example:

\[
\boxed{
R_e:
\Phi_e^A
\neq
B
}
\]

may ensure that a particular oracle computation does not compute set \(B\).

The construction must satisfy all requirements eventually.

%%%%%%%%%%%%%%%%%%%%%%%%%%%%%%%%%%%%%%%%%%%%%%%%%%%%%%%%%%%%
\subsection{Injury}
\label{subsec:priority-injury}
%%%%%%%%%%%%%%%%%%%%%%%%%%%%%%%%%%%%%%%%%%%%%%%%%%%%%%%%%%%%

A higher-priority requirement may force a change that destroys progress
made by a lower-priority requirement.

This is called injury.

The lower-priority strategy must then restart or repair itself.

The construction succeeds when every requirement is injured only in a
controlled manner.

%%%%%%%%%%%%%%%%%%%%%%%%%%%%%%%%%%%%%%%%%%%%%%%%%%%%%%%%%%%%
\subsection{Finite-Injury Priority Arguments}
\label{subsec:finite-injury}
%%%%%%%%%%%%%%%%%%%%%%%%%%%%%%%%%%%%%%%%%%%%%%%%%%%%%%%%%%%%

In a finite-injury argument, every requirement is injured only finitely
many times.

Eventually each requirement reaches a stable stage and permanently
satisfies its goal.

The proof often proceeds inductively through the priority ordering.

%%%%%%%%%%%%%%%%%%%%%%%%%%%%%%%%%%%%%%%%%%%%%%%%%%%%%%%%%%%%
\subsection{Priority Method Intuition}
\label{subsec:priority-method-intuition}
%%%%%%%%%%%%%%%%%%%%%%%%%%%%%%%%%%%%%%%%%%%%%%%%%%%%%%%%%%%%

Priority arguments are similar to managing competing constraints:

\[
\boxed{
R_0
>
R_1
>
R_2
>
\cdots.
}
\]

Requirement \(R_0\) gets first choice.

Requirement \(R_1\) acts only in ways compatible with \(R_0\), and so on.

If higher-priority actions later change, lower-priority work may need
revision.

%%%%%%%%%%%%%%%%%%%%%%%%%%%%%%%%%%%%%%%%%%%%%%%%%%%%%%%%%%%%
\subsection{Arithmetic Hierarchy}
\label{subsec:recursion-arithmetic-hierarchy}
%%%%%%%%%%%%%%%%%%%%%%%%%%%%%%%%%%%%%%%%%%%%%%%%%%%%%%%%%%%%

The arithmetic hierarchy classifies subsets of

\[
\mathbb{N}
\]

according to the complexity of unbounded quantifiers in their arithmetic
definitions.

The principal classes are

\[
\boxed{
\Sigma^0_n,
\qquad
\Pi^0_n,
\qquad
\Delta^0_n.
}
\]

The superscript \(0\) indicates first-order arithmetic quantification over
natural numbers.

%%%%%%%%%%%%%%%%%%%%%%%%%%%%%%%%%%%%%%%%%%%%%%%%%%%%%%%%%%%%
\subsection{Bounded Quantifiers}
\label{subsec:bounded-quantifiers-arithmetic-hierarchy}
%%%%%%%%%%%%%%%%%%%%%%%%%%%%%%%%%%%%%%%%%%%%%%%%%%%%%%%%%%%%

Quantifiers such as

\[
\boxed{
\exists x<t
}
\]

or

\[
\boxed{
\forall x<t
}
\]

are bounded.

Bounded quantifiers do not raise the arithmetic-hierarchy level in the same
way unrestricted quantifiers do.

The hierarchy measures alternations of unbounded number quantifiers.

%%%%%%%%%%%%%%%%%%%%%%%%%%%%%%%%%%%%%%%%%%%%%%%%%%%%%%%%%%%%
\subsection{\(\Sigma^0_1\) Sets}
\label{subsec:sigma-zero-one}
%%%%%%%%%%%%%%%%%%%%%%%%%%%%%%%%%%%%%%%%%%%%%%%%%%%%%%%%%%%%

A set \(A\) is \(\Sigma^0_1\) if membership can be expressed in form

\[
\boxed{
x\in A
\iff
\exists y\,
R(x,y),
}
\]

where \(R\) is computable, allowing equivalent standard normal forms.

These are exactly the computably enumerable sets.

Thus

\[
\boxed{
\Sigma^0_1
=
\text{c.e.}
}
\]

at the level of subsets of \(\mathbb{N}\).

%%%%%%%%%%%%%%%%%%%%%%%%%%%%%%%%%%%%%%%%%%%%%%%%%%%%%%%%%%%%
\subsection{\(\Pi^0_1\) Sets}
\label{subsec:pi-zero-one}
%%%%%%%%%%%%%%%%%%%%%%%%%%%%%%%%%%%%%%%%%%%%%%%%%%%%%%%%%%%%

A set \(A\) is \(\Pi^0_1\) when membership can be written

\[
\boxed{
x\in A
\iff
\forall y\,
R(x,y),
}
\]

with computable \(R\).

These correspond to complements of c.e.\ sets.

Therefore

\[
\boxed{
\Pi^0_1
=
\text{co-c.e.}
}
\]

in the standard arithmetic setting.

%%%%%%%%%%%%%%%%%%%%%%%%%%%%%%%%%%%%%%%%%%%%%%%%%%%%%%%%%%%%
\subsection{\(\Delta^0_1\) Sets}
\label{subsec:delta-zero-one}
%%%%%%%%%%%%%%%%%%%%%%%%%%%%%%%%%%%%%%%%%%%%%%%%%%%%%%%%%%%%

The class

\[
\Delta^0_1
\]

consists of sets that are both

\[
\Sigma^0_1
\]

and

\[
\Pi^0_1.
\]

Therefore,

\[
\boxed{
\Delta^0_1
=
\text{computable sets}.
}
\]

This recovers the theorem that c.e.\ plus co-c.e.\ implies computable.

%%%%%%%%%%%%%%%%%%%%%%%%%%%%%%%%%%%%%%%%%%%%%%%%%%%%%%%%%%%%
\subsection{Higher Arithmetic Levels}
\label{subsec:higher-arithmetic-levels}
%%%%%%%%%%%%%%%%%%%%%%%%%%%%%%%%%%%%%%%%%%%%%%%%%%%%%%%%%%%%

A typical \(\Sigma^0_2\) definition has form

\[
\boxed{
x\in A
\iff
\exists y
\forall z
\,
R(x,y,z),
}
\]

with computable \(R\).

A typical \(\Pi^0_2\) definition reverses the leading quantifier:

\[
\boxed{
x\in A
\iff
\forall y
\exists z
\,
R(x,y,z).
}
\]

Each additional unbounded alternation can increase definitional complexity.

%%%%%%%%%%%%%%%%%%%%%%%%%%%%%%%%%%%%%%%%%%%%%%%%%%%%%%%%%%%%
\subsection{\(\Sigma^0_n\) and \(\Pi^0_n\)}
\label{subsec:sigma-pi-n}
%%%%%%%%%%%%%%%%%%%%%%%%%%%%%%%%%%%%%%%%%%%%%%%%%%%%%%%%%%%%

Informally:

\[
\boxed{
\Sigma^0_n
}
\]

begins with an existential block of unbounded quantifiers and contains at
most \(n\) alternating blocks.

Similarly,

\[
\boxed{
\Pi^0_n
}
\]

begins with a universal block.

Precise definitions are usually given recursively or through normal-form
formulas.

%%%%%%%%%%%%%%%%%%%%%%%%%%%%%%%%%%%%%%%%%%%%%%%%%%%%%%%%%%%%
\subsection{\(\Delta^0_n\)}
\label{subsec:delta-zero-n}
%%%%%%%%%%%%%%%%%%%%%%%%%%%%%%%%%%%%%%%%%%%%%%%%%%%%%%%%%%%%

For \(n\geq1\),

\[
\boxed{
\Delta^0_n
=
\Sigma^0_n
\cap
\Pi^0_n
}
\]

under the standard hierarchy conventions.

These sets are definable at the same level from both existential-first and
universal-first viewpoints.

%%%%%%%%%%%%%%%%%%%%%%%%%%%%%%%%%%%%%%%%%%%%%%%%%%%%%%%%%%%%
\subsection{Post's Theorem}
\label{subsec:posts-theorem}
%%%%%%%%%%%%%%%%%%%%%%%%%%%%%%%%%%%%%%%%%%%%%%%%%%%%%%%%%%%%

Post's theorem connects the arithmetic hierarchy with Turing jumps.

A central form states that

\[
\boxed{
\Sigma^0_{n+1}
}
\]

corresponds to sets computably enumerable relative to

\[
\boxed{
\mathbf{0}^{(n)},
}
\]

while

\[
\boxed{
\Delta^0_{n+1}
}
\]

corresponds to sets computable relative to

\[
\boxed{
\mathbf{0}^{(n)}.
}
\]

Thus logical quantifier complexity matches levels of oracle computation.

%%%%%%%%%%%%%%%%%%%%%%%%%%%%%%%%%%%%%%%%%%%%%%%%%%%%%%%%%%%%
\subsection{First Levels Under Post's Theorem}
\label{subsec:posts-theorem-first-levels}
%%%%%%%%%%%%%%%%%%%%%%%%%%%%%%%%%%%%%%%%%%%%%%%%%%%%%%%%%%%%

At the bottom,

\[
\boxed{
\Delta^0_1
=
\text{computable}.
}
\]

Next,

\[
\boxed{
\Sigma^0_1
=
\text{c.e.},
}
\]

and

\[
\boxed{
\Pi^0_1
=
\text{co-c.e.}
}
\]

At the next level,

\[
\Delta^0_2
\]

consists of sets computable relative to

\[
\boxed{
\mathbf{0}'.
}
\]

%%%%%%%%%%%%%%%%%%%%%%%%%%%%%%%%%%%%%%%%%%%%%%%%%%%%%%%%%%%%
\subsection{Complete Sets in the Arithmetic Hierarchy}
\label{subsec:arithmetic-hierarchy-complete-sets}
%%%%%%%%%%%%%%%%%%%%%%%%%%%%%%%%%%%%%%%%%%%%%%%%%%%%%%%%%%%%

A set may be complete for a hierarchy level under a specified reducibility.

For example, the halting set is

\[
\boxed{
\Sigma^0_1\text{-complete}
}
\]

under standard many-one reducibility.

Complete sets represent maximal difficulty within their corresponding
definability class under the chosen reduction.

%%%%%%%%%%%%%%%%%%%%%%%%%%%%%%%%%%%%%%%%%%%%%%%%%%%%%%%%%%%%
\subsection{Hierarchy Properness}
\label{subsec:arithmetic-hierarchy-properness}
%%%%%%%%%%%%%%%%%%%%%%%%%%%%%%%%%%%%%%%%%%%%%%%%%%%%%%%%%%%%

The arithmetic hierarchy does not collapse at finite levels.

There are sets lying genuinely higher in the hierarchy.

Conceptually,

\[
\boxed{
\Sigma^0_n
\cup
\Pi^0_n
\subsetneq
\Delta^0_{n+1}
}
\]

in the standard hierarchy.

The jump operator provides a mechanism for proving strictness.

%%%%%%%%%%%%%%%%%%%%%%%%%%%%%%%%%%%%%%%%%%%%%%%%%%%%%%%%%%%%
\subsection{Computable Approximations}
\label{subsec:computable-approximations}
%%%%%%%%%%%%%%%%%%%%%%%%%%%%%%%%%%%%%%%%%%%%%%%%%%%%%%%%%%%%

A noncomputable set may sometimes be approached through a sequence of
computable guesses.

Let

\[
\boxed{
A_s(x)
\in
\{0,1\}
}
\]

be the stage-\(s\) approximation to membership of \(x\).

If

\[
\boxed{
\lim_{s\to\infty}
A_s(x)
=
\chi_A(x),
}
\]

then the guesses eventually stabilize to the correct answer for each
individual \(x\).

%%%%%%%%%%%%%%%%%%%%%%%%%%%%%%%%%%%%%%%%%%%%%%%%%%%%%%%%%%%%
\subsection{Limit Computability}
\label{subsec:limit-computability}
%%%%%%%%%%%%%%%%%%%%%%%%%%%%%%%%%%%%%%%%%%%%%%%%%%%%%%%%%%%%

A function \(f\) is limit computable if there exists a computable function

\[
g(x,s)
\]

such that

\[
\boxed{
f(x)
=
\lim_{s\to\infty}
g(x,s),
}
\]

where the limit is understood as eventual stabilization.

The computation is allowed to revise its guess finitely many times for each
input.

%%%%%%%%%%%%%%%%%%%%%%%%%%%%%%%%%%%%%%%%%%%%%%%%%%%%%%%%%%%%
\subsection{Limit Lemma}
\label{subsec:limit-lemma}
%%%%%%%%%%%%%%%%%%%%%%%%%%%%%%%%%%%%%%%%%%%%%%%%%%%%%%%%%%%%

A central limit lemma states, in one standard form, that sets computable
from the halting problem are exactly those having suitable computable
convergent approximations.

Conceptually,

\[
\boxed{
A\leq_T\mathbf{0}'
}
\]

if and only if

\[
\boxed{
A
\text{ has a computable approximation converging pointwise}.
}
\]

This connects oracle computation with changing finite-stage guesses.

%%%%%%%%%%%%%%%%%%%%%%%%%%%%%%%%%%%%%%%%%%%%%%%%%%%%%%%%%%%%
\subsection{Worked Example: Approximation That Changes}
\label{subsec:worked-computable-approximation}
%%%%%%%%%%%%%%%%%%%%%%%%%%%%%%%%%%%%%%%%%%%%%%%%%%%%%%%%%%%%

\begin{workedexamplebox}{Approximating a c.e.\ Set}

Let \(A\) be c.e.\ and let

\[
A_s
\]

contain the elements enumerated into \(A\) by stage \(s\).

Define

\[
\chi_s(x)
=
\begin{cases}
1,&x\in A_s,\\
0,&x\notin A_s.
\end{cases}
\]

If \(x\in A\), eventually \(x\) is enumerated and all later guesses equal
\(1\).

If \(x\notin A\), every guess remains \(0\).

Thus

\begin{answerbox}
\[
\boxed{
\lim_{s\to\infty}
\chi_s(x)
=
\chi_A(x).
}
\]
\end{answerbox}

The approximation is computable even when \(A\) itself is not decidable.

\end{workedexamplebox}

%%%%%%%%%%%%%%%%%%%%%%%%%%%%%%%%%%%%%%%%%%%%%%%%%%%%%%%%%%%%
\subsection{Degree Structures Are Partial Orders}
\label{subsec:degree-partial-order}
%%%%%%%%%%%%%%%%%%%%%%%%%%%%%%%%%%%%%%%%%%%%%%%%%%%%%%%%%%%%

Turing degrees are not generally linearly ordered.

There may exist degrees

\[
\mathbf{a}
\]

and

\[
\mathbf{b}
\]

such that neither

\[
\mathbf{a}\leq_T\mathbf{b}
\]

nor

\[
\mathbf{b}\leq_T\mathbf{a}.
\]

Thus noncomputability has a branching structure rather than one simple
numerical scale.

%%%%%%%%%%%%%%%%%%%%%%%%%%%%%%%%%%%%%%%%%%%%%%%%%%%%%%%%%%%%
\subsection{Joins of Degrees}
\label{subsec:join-turing-degrees}
%%%%%%%%%%%%%%%%%%%%%%%%%%%%%%%%%%%%%%%%%%%%%%%%%%%%%%%%%%%%

Given sets \(A\) and \(B\), combine their information into a set such as

\[
\boxed{
A\oplus B
=
\{2n:n\in A\}
\cup
\{2n+1:n\in B\}.
}
\]

Then \(A\oplus B\) computes both \(A\) and \(B\).

Its degree acts as a least upper bound of their Turing degrees.

%%%%%%%%%%%%%%%%%%%%%%%%%%%%%%%%%%%%%%%%%%%%%%%%%%%%%%%%%%%%
\subsection{Worked Example: Join Coding}
\label{subsec:worked-degree-join}
%%%%%%%%%%%%%%%%%%%%%%%%%%%%%%%%%%%%%%%%%%%%%%%%%%%%%%%%%%%%

\begin{workedexamplebox}{Recovering Two Sets From Their Join}

Let

\[
C=A\oplus B.
\]

To determine whether

\[
n\in A,
\]

ask whether

\[
2n\in C.
\]

To determine whether

\[
n\in B,
\]

ask whether

\[
2n+1\in C.
\]

Therefore,

\begin{answerbox}
\[
\boxed{
A\leq_TC
\qquad\text{and}\qquad
B\leq_TC.
}
\]
\end{answerbox}

\end{workedexamplebox}

%%%%%%%%%%%%%%%%%%%%%%%%%%%%%%%%%%%%%%%%%%%%%%%%%%%%%%%%%%%%
\subsection{Enumeration Degrees Preview}
\label{subsec:enumeration-degrees-preview}
%%%%%%%%%%%%%%%%%%%%%%%%%%%%%%%%%%%%%%%%%%%%%%%%%%%%%%%%%%%%

Turing reducibility assumes oracle membership questions receive both
positive and negative answers.

Enumeration reducibility studies information supplied only through positive
enumeration.

This produces another degree structure suited to c.e.-style information.

%%%%%%%%%%%%%%%%%%%%%%%%%%%%%%%%%%%%%%%%%%%%%%%%%%%%%%%%%%%%
\subsection{Uniformity}
\label{subsec:recursion-theory-uniformity}
%%%%%%%%%%%%%%%%%%%%%%%%%%%%%%%%%%%%%%%%%%%%%%%%%%%%%%%%%%%%

Recursion theory often distinguishes between:

\[
\boxed{
\text{an object existing for every input}
}
\]

and

\[
\boxed{
\text{one algorithm uniformly producing those objects}.
}
\]

The \(s\)-\(m\)-\(n\) theorem is a major uniformity result.

Uniformity is essential when transforming mathematical existence statements
into effective constructions.

%%%%%%%%%%%%%%%%%%%%%%%%%%%%%%%%%%%%%%%%%%%%%%%%%%%%%%%%%%%%
\subsection{Effective Versus Noneffective Existence}
\label{subsec:effective-noneffective-existence}
%%%%%%%%%%%%%%%%%%%%%%%%%%%%%%%%%%%%%%%%%%%%%%%%%%%%%%%%%%%%

A set-theoretic proof may show:

\[
\boxed{
\forall x\exists y\,P(x,y).
}
\]

This does not automatically produce a computable function \(f\) satisfying

\[
\boxed{
P(x,f(x))
}
\]

for every \(x\).

Recursion theory studies when witnesses can be selected effectively.

%%%%%%%%%%%%%%%%%%%%%%%%%%%%%%%%%%%%%%%%%%%%%%%%%%%%%%%%%%%%
\subsection{Recursion Theory and Arithmetic}
\label{subsec:recursion-theory-arithmetic}
%%%%%%%%%%%%%%%%%%%%%%%%%%%%%%%%%%%%%%%%%%%%%%%%%%%%%%%%%%%%

Because finite computation can be encoded arithmetically, recursion theory
and arithmetic are deeply connected.

Computational statements become statements about natural numbers.

For example,

\[
\boxed{
\varphi_e(x)\downarrow
}
\]

can be represented by the existence of a coded finite computation history.

Thus computability properties acquire arithmetic definitions.

%%%%%%%%%%%%%%%%%%%%%%%%%%%%%%%%%%%%%%%%%%%%%%%%%%%%%%%%%%%%
\subsection{Recursion Theory and Incompleteness}
\label{subsec:recursion-theory-incompleteness}
%%%%%%%%%%%%%%%%%%%%%%%%%%%%%%%%%%%%%%%%%%%%%%%%%%%%%%%%%%%%

The theorem set of an effective formal theory is computably enumerable.

If a sufficiently strong theory were both complete and effectively
axiomatized under suitable consistency hypotheses, one could often obtain
a decision procedure for arithmetic truth that is known not to exist.

This provides one computational perspective on incompleteness.

%%%%%%%%%%%%%%%%%%%%%%%%%%%%%%%%%%%%%%%%%%%%%%%%%%%%%%%%%%%%
\subsection{Recursion Theory and Definability}
\label{subsec:recursion-definability}
%%%%%%%%%%%%%%%%%%%%%%%%%%%%%%%%%%%%%%%%%%%%%%%%%%%%%%%%%%%%

The arithmetic hierarchy translates computational complexity of
information into logical definability.

For example,

\[
\boxed{
\text{c.e.}
\leftrightarrow
\Sigma^0_1.
}
\]

Higher quantifier alternations correspond to higher oracle jumps.

Thus:

\[
\boxed{
\text{computation}
}
\]

and

\[
\boxed{
\text{logical definition}
}
\]

form parallel hierarchies.

%%%%%%%%%%%%%%%%%%%%%%%%%%%%%%%%%%%%%%%%%%%%%%%%%%%%%%%%%%%%
\subsection{Recursion Theory Workflow}
\label{subsec:recursion-theory-workflow}
%%%%%%%%%%%%%%%%%%%%%%%%%%%%%%%%%%%%%%%%%%%%%%%%%%%%%%%%%%%%

When studying a recursion-theoretic problem:

\begin{enumerate}
    \item encode all finite objects by natural numbers;
    \item determine whether functions are total or partial;
    \item identify whether primitive recursion is sufficient;
    \item use minimization when unbounded search is required;
    \item express computable functions through effective indices;
    \item use universal computation to manipulate program descriptions;
    \item apply the \(s\)-\(m\)-\(n\) theorem when parameters must be
          hard-coded;
    \item apply the recursion theorem when self-reference or program
          fixed points are required;
    \item identify whether a set is computable, c.e., or co-c.e.;
    \item express c.e.\ sets as \(W_e\) when working uniformly with
          enumerations;
    \item choose the appropriate reducibility;
    \item compare sets through Turing degrees when oracle power matters;
    \item apply the jump to move to a stronger information level;
    \item identify the arithmetic-hierarchy complexity of a definition;
    \item count unbounded quantifier alternations carefully;
    \item apply Post's theorem to connect definability with oracle
          computation;
    \item use computable approximations when analyzing
          \(\Delta^0_2\)-type behavior;
    \item use priority constructions when many incompatible requirements
          must be satisfied;
    \item distinguish existence of an intermediate object from a uniform
          method for producing it;
    \item track all oracle dependencies explicitly.
\end{enumerate}

%%%%%%%%%%%%%%%%%%%%%%%%%%%%%%%%%%%%%%%%%%%%%%%%%%%%%%%%%%%%
\subsection{Common Errors}
\label{subsec:recursion-theory-common-errors}
%%%%%%%%%%%%%%%%%%%%%%%%%%%%%%%%%%%%%%%%%%%%%%%%%%%%%%%%%%%%

\subsubsection{Confusing Recursion Theory With Recursive Programming Only}

Recursion theory studies all computability, not merely programs written
with recursive function calls.

Its subject includes reducibility, degrees, jumps, definability, and
noncomputability.

%%%%%%%%%%%%%%%%%%%%%%%%%%%%%%%%%%%%%%%%%%%%%%%%%%%%%%%%%%%%
\subsubsection{Assuming Primitive Recursive Means Partially Defined}

Primitive recursive functions are total.

Unbounded minimization is what naturally introduces partiality.

%%%%%%%%%%%%%%%%%%%%%%%%%%%%%%%%%%%%%%%%%%%%%%%%%%%%%%%%%%%%
\subsubsection{Assuming Every Total Computable Function Is Primitive Recursive}

The Ackermann function provides a standard counterexample.

Thus

\[
\boxed{
\text{primitive recursive}
\subsetneq
\text{total computable}.
}
\]

%%%%%%%%%%%%%%%%%%%%%%%%%%%%%%%%%%%%%%%%%%%%%%%%%%%%%%%%%%%%
\subsubsection{Confusing Primitive Recursion With General Recursion}

Primitive recursion uses a controlled recursion over a natural-number
argument.

General computability allows more powerful mechanisms such as unbounded
search.

%%%%%%%%%%%%%%%%%%%%%%%%%%%%%%%%%%%%%%%%%%%%%%%%%%%%%%%%%%%%
\subsubsection{Assuming Minimization Always Terminates}

The search

\[
\mu y\,[R(x,y)]
\]

may continue forever if no witness exists.

That is precisely why minimization naturally produces partial functions.

%%%%%%%%%%%%%%%%%%%%%%%%%%%%%%%%%%%%%%%%%%%%%%%%%%%%%%%%%%%%
\subsubsection{Assuming a Function Has One Unique Program Index}

A computable function usually has infinitely many different program
descriptions and therefore many indices.

Indices belong to descriptions, not uniquely to functions.

%%%%%%%%%%%%%%%%%%%%%%%%%%%%%%%%%%%%%%%%%%%%%%%%%%%%%%%%%%%%
\subsubsection{Confusing Program Equality With Index Equality}

If

\[
e\neq d,
\]

the programs may still compute the same partial function:

\[
\varphi_e=\varphi_d.
\]

%%%%%%%%%%%%%%%%%%%%%%%%%%%%%%%%%%%%%%%%%%%%%%%%%%%%%%%%%%%%
\subsubsection{Treating the \(s\)-\(m\)-\(n\) Theorem as Ordinary Substitution}

The theorem is stronger than a symbolic substitution observation.

It asserts an effective transformation of program indices.

%%%%%%%%%%%%%%%%%%%%%%%%%%%%%%%%%%%%%%%%%%%%%%%%%%%%%%%%%%%%
\subsubsection{Interpreting the Recursion Theorem as Ordinary Numeric Recursion}

Kleene's recursion theorem is a fixed-point theorem for program behavior.

It is different from defining a sequence by

\[
a_{n+1}=F(a_n).
\]

%%%%%%%%%%%%%%%%%%%%%%%%%%%%%%%%%%%%%%%%%%%%%%%%%%%%%%%%%%%%
\subsubsection{Assuming c.e.\ Means Computable}

A c.e.\ set permits positive membership to be discovered eventually.

It may still have no decision algorithm.

%%%%%%%%%%%%%%%%%%%%%%%%%%%%%%%%%%%%%%%%%%%%%%%%%%%%%%%%%%%%
\subsubsection{Assuming Nonmembership in a c.e.\ Set Can Be Confirmed}

If an element has not yet been enumerated, it may appear later.

Failure to appear by stage \(s\) does not prove permanent nonmembership.

%%%%%%%%%%%%%%%%%%%%%%%%%%%%%%%%%%%%%%%%%%%%%%%%%%%%%%%%%%%%
\subsubsection{Confusing co-c.e.\ With Complement of Computable}

The complement of a computable set is computable.

A co-c.e.\ set need not itself be computable.

It merely has a c.e.\ complement.

%%%%%%%%%%%%%%%%%%%%%%%%%%%%%%%%%%%%%%%%%%%%%%%%%%%%%%%%%%%%
\subsubsection{Assuming Every Noncomputable c.e.\ Set Is Complete}

Simple and intermediate c.e.\ sets show this is false.

Noncomputability has many degrees.

%%%%%%%%%%%%%%%%%%%%%%%%%%%%%%%%%%%%%%%%%%%%%%%%%%%%%%%%%%%%
\subsubsection{Confusing Many-One and Turing Reducibility}

Many-one reducibility permits a single effective transformation.

Turing reducibility permits repeated adaptive oracle queries.

Thus

\[
\boxed{
\leq_m
}
\]

is more restrictive than

\[
\boxed{
\leq_T.
}
\]

%%%%%%%%%%%%%%%%%%%%%%%%%%%%%%%%%%%%%%%%%%%%%%%%%%%%%%%%%%%%
\subsubsection{Assuming Turing Degrees Form a Linear Order}

Degrees are only partially ordered.

Two sets may have incomparable computational information.

%%%%%%%%%%%%%%%%%%%%%%%%%%%%%%%%%%%%%%%%%%%%%%%%%%%%%%%%%%%%
\subsubsection{Assuming the Halting Degree Is the Highest Degree}

The halting degree

\[
\mathbf{0}'
\]

is above computable sets, but one may form

\[
\boxed{
\mathbf{0}'',
\mathbf{0}''',
\ldots.
}
\]

There are also many other degrees not lying on this simple jump sequence.

%%%%%%%%%%%%%%%%%%%%%%%%%%%%%%%%%%%%%%%%%%%%%%%%%%%%%%%%%%%%
\subsubsection{Assuming an Oracle Solves Its Own Halting Problem}

An oracle for \(A\) does not compute \(A'\).

Indeed,

\[
\boxed{
A<_TA'.
}
\]

%%%%%%%%%%%%%%%%%%%%%%%%%%%%%%%%%%%%%%%%%%%%%%%%%%%%%%%%%%%%
\subsubsection{Confusing Post's Problem With Post's Theorem}

Post's problem asks about intermediate c.e.\ degrees.

Post's theorem relates the arithmetic hierarchy to iterated Turing jumps.

They are different results.

%%%%%%%%%%%%%%%%%%%%%%%%%%%%%%%%%%%%%%%%%%%%%%%%%%%%%%%%%%%%
\subsubsection{Counting Bounded Quantifiers as New Arithmetic-Hierarchy Levels}

The arithmetic hierarchy is determined by unbounded quantifier
alternations.

Bounded quantifiers do not raise the level in the standard classification.

%%%%%%%%%%%%%%%%%%%%%%%%%%%%%%%%%%%%%%%%%%%%%%%%%%%%%%%%%%%%
\subsubsection{Assuming \(\Sigma^0_n\) and \(\Pi^0_n\) Are Identical}

At nontrivial finite levels, these classes are distinct.

Their intersection gives the corresponding \(\Delta\)-class in the standard
hierarchy.

%%%%%%%%%%%%%%%%%%%%%%%%%%%%%%%%%%%%%%%%%%%%%%%%%%%%%%%%%%%%
\subsubsection{Confusing Logical Hierarchy With Runtime Complexity}

The arithmetic hierarchy measures definability and oracle complexity.

It is not the same as polynomial time, exponential time, or other runtime
complexity classes.

%%%%%%%%%%%%%%%%%%%%%%%%%%%%%%%%%%%%%%%%%%%%%%%%%%%%%%%%%%%%
\subsubsection{Assuming a Computable Approximation Must Always Be Correct}

A computable approximation may make temporary mistakes.

The requirement is eventual stabilization to the correct value.

%%%%%%%%%%%%%%%%%%%%%%%%%%%%%%%%%%%%%%%%%%%%%%%%%%%%%%%%%%%%
\subsubsection{Assuming Priority Requirements Never Interfere}

Priority constructions are specifically designed for interacting
requirements.

Higher-priority strategies may injure lower-priority ones.

%%%%%%%%%%%%%%%%%%%%%%%%%%%%%%%%%%%%%%%%%%%%%%%%%%%%%%%%%%%%
\subsubsection{Assuming Finite Injury Means Only Finitely Many Injuries Overall}

In a countable construction, there may be infinitely many injuries across
all requirements.

Finite injury means each individual requirement is injured only finitely
many times.

%%%%%%%%%%%%%%%%%%%%%%%%%%%%%%%%%%%%%%%%%%%%%%%%%%%%%%%%%%%%
\subsection{Applications}
\label{subsec:recursion-theory-applications}
%%%%%%%%%%%%%%%%%%%%%%%%%%%%%%%%%%%%%%%%%%%%%%%%%%%%%%%%%%%%

\begin{applicationbox}

\textbf{Foundations of Computation.}
Recursive-function definitions provide an arithmetic characterization of
the same computable functions described by Turing machines.

\medskip

\textbf{Mathematical Logic.}
Computably enumerable theories, Gödel coding, self-reference, and
arithmetical definability are naturally analyzed using recursion theory.

\medskip

\textbf{Proof Theory.}
The effective enumeration of formal proofs connects theoremhood with c.e.\
sets and undecidability.

\medskip

\textbf{Arithmetic.}
The arithmetic hierarchy classifies number-theoretic predicates according
to quantifier complexity and computational strength.

\medskip

\textbf{Program Semantics.}
The \(s\)-\(m\)-\(n\) and recursion theorems formalize parameterized and
self-referential program behavior.

\medskip

\textbf{Relative Computability.}
Oracle computation and Turing degrees measure how different sources of
noncomputable information compare.

\medskip

\textbf{Definability Theory.}
Post's theorem connects logical descriptions with iterated computational
jumps.

\medskip

\textbf{Classical Recursion Theory.}
Priority methods, simple sets, intermediate degrees, and degree structures
form major areas of mathematical research.

\end{applicationbox}

%%%%%%%%%%%%%%%%%%%%%%%%%%%%%%%%%%%%%%%%%%%%%%%%%%%%%%%%%%%%
\subsection{Exercises}
\label{subsec:recursion-theory-exercises}
%%%%%%%%%%%%%%%%%%%%%%%%%%%%%%%%%%%%%%%%%%%%%%%%%%%%%%%%%%%%

\begin{exercise}[1]
Define a primitive recursive function.
\end{exercise}

\begin{exercise}[1]
List the initial primitive recursive functions.
\end{exercise}

\begin{exercise}[1]
Define composition.
\end{exercise}

\begin{exercise}[1]
Define primitive recursion.
\end{exercise}

\begin{exercise}[1]
Define unbounded minimization.
\end{exercise}

\begin{exercise}[1]
Define a partial recursive function.
\end{exercise}

\begin{exercise}[1]
Define a recursive function in the total-function sense.
\end{exercise}

\begin{exercise}[1]
Define a computably enumerable set.
\end{exercise}

\begin{exercise}[1]
Define Turing reducibility.
\end{exercise}

\begin{exercise}[1]
Define a Turing degree.
\end{exercise}

\begin{exercise}[2]
Show that addition is primitive recursive.
\end{exercise}

\begin{exercise}[2]
Show that multiplication is primitive recursive using addition.
\end{exercise}

\begin{exercise}[2]
Show that exponentiation is primitive recursive using multiplication.
\end{exercise}

\begin{exercise}[2]
Show that factorial is primitive recursive.
\end{exercise}

\begin{exercise}[3]
Construct primitive recursive predecessor.
\end{exercise}

\begin{exercise}[3]
Construct truncated subtraction.
\end{exercise}

\begin{exercise}[3]
Show that equality is a primitive recursive predicate.
\end{exercise}

\begin{exercise}[3]
Show that

\[
x\leq y
\]

is primitive recursive.
\end{exercise}

\begin{exercise}[3]
Explain why bounded search preserves primitive recursiveness.
\end{exercise}

\begin{exercise}[3]
Explain why bounded existential quantification preserves primitive
recursiveness.
\end{exercise}

\begin{exercise}[3]
Explain why every primitive recursive function is total.
\end{exercise}

\begin{exercise}[2]
Describe the Ackermann function conceptually.
\end{exercise}

\begin{exercise}[3]
Explain why Ackermann-type growth demonstrates that primitive recursive
functions do not include every total computable function.
\end{exercise}

\begin{exercise}[2]
Explain why unbounded minimization may fail to terminate.
\end{exercise}

\begin{exercise}[3]
Give an example of a partial function produced by unbounded search.
\end{exercise}

\begin{exercise}[3]
Explain why partial recursive and partial Turing-computable functions define
the same class.
\end{exercise}

\begin{exercise}[2]
Explain what an index \(e\) represents in

\[
\varphi_e.
\]
\end{exercise}

\begin{exercise}[2]
Explain why indices are not unique.
\end{exercise}

\begin{exercise}[3]
Define a universal partial computable function conceptually.
\end{exercise}

\begin{exercise}[3]
Explain the \(s\)-\(m\)-\(n\) theorem conceptually.
\end{exercise}

\begin{exercise}[3]
Use parameterization to specialize a two-variable computable function at one
fixed first argument.
\end{exercise}

\begin{exercise}[3]
Explain why the \(s\)-\(m\)-\(n\) theorem is an effective uniformity result.
\end{exercise}

\begin{exercise}[3]
Explain Kleene normal form conceptually.
\end{exercise}

\begin{exercise}[3]
Explain why finite computation histories are primitive recursively
verifiable.
\end{exercise}

\begin{exercise}[3]
State Kleene's recursion theorem.
\end{exercise}

\begin{exercise}[3]
Explain why the recursion theorem is a fixed-point theorem.
\end{exercise}

\begin{exercise}[3]
Explain how it permits self-reference without direct access to source code.
\end{exercise}

\begin{exercise}[3]
Compare the recursion theorem with Gödel's diagonal lemma.
\end{exercise}

\begin{exercise}[2]
Define

\[
W_e.
\]
\end{exercise}

\begin{exercise}[3]
Explain finite-stage approximations

\[
W_{e,s}.
\]
\end{exercise}

\begin{exercise}[3]
Show that if \(A\) and its complement are c.e., then \(A\) is computable.
\end{exercise}

\begin{exercise}[2]
Define a co-c.e.\ set.
\end{exercise}

\begin{exercise}[3]
Explain why the diagonal halting set \(K\) is c.e.
\end{exercise}

\begin{exercise}[3]
Explain why \(K\) is not computable.
\end{exercise}

\begin{exercise}[3]
Define c.e.\ many-one completeness.
\end{exercise}

\begin{exercise}[3]
Explain why a complete c.e.\ set is maximally difficult among c.e.\ sets
under the chosen reduction.
\end{exercise}

\begin{exercise}[3]
Define an index set.
\end{exercise}

\begin{exercise}[3]
Explain why semantic program properties naturally give index sets.
\end{exercise}

\begin{exercise}[3]
Define a productive set conceptually.
\end{exercise}

\begin{exercise}[3]
Define a creative set conceptually.
\end{exercise}

\begin{exercise}[3]
Explain computable inseparability.
\end{exercise}

\begin{exercise}[3]
Define an immune set.
\end{exercise}

\begin{exercise}[3]
Define a simple set.
\end{exercise}

\begin{exercise}[3]
Explain why a simple set demonstrates that noncomputable c.e.\ sets need not
be complete.
\end{exercise}

\begin{exercise}[2]
Define one-one reducibility.
\end{exercise}

\begin{exercise}[3]
Show conceptually that

\[
A\leq_1B
\Longrightarrow
A\leq_mB.
\]
\end{exercise}

\begin{exercise}[2]
Define truth-table reducibility conceptually.
\end{exercise}

\begin{exercise}[3]
Compare adaptive and nonadaptive oracle queries.
\end{exercise}

\begin{exercise}[3]
Explain why

\[
A\leq_mB
\Longrightarrow
A\leq_TB.
\]
\end{exercise}

\begin{exercise}[2]
Define Turing equivalence.
\end{exercise}

\begin{exercise}[2]
Define the computable degree

\[
\mathbf{0}.
\]
\end{exercise}

\begin{exercise}[2]
Define the halting degree

\[
\mathbf{0}'.
\]
\end{exercise}

\begin{exercise}[3]
Explain why

\[
\mathbf{0}
<
\mathbf{0}'.
\]
\end{exercise}

\begin{exercise}[3]
Define the Turing jump of a set \(A\).
\end{exercise}

\begin{exercise}[3]
Explain why

\[
A<_TA'.
\]
\end{exercise}

\begin{exercise}[3]
Describe the iterated jumps

\[
\mathbf{0}',
\mathbf{0}'',
\mathbf{0}'''.
\]
\end{exercise}

\begin{exercise}[3]
Explain relativization.
\end{exercise}

\begin{exercise}[3]
Explain why oracle \(A\) does not decide its own relative halting problem.
\end{exercise}

\begin{exercise}[2]
State Post's problem.
\end{exercise}

\begin{exercise}[3]
Explain what an intermediate c.e.\ degree is.
\end{exercise}

\begin{exercise}[3]
Explain why the existence of intermediate degrees reveals structure inside
the c.e.\ sets.
\end{exercise}

\begin{exercise}[2]
Describe the priority method conceptually.
\end{exercise}

\begin{exercise}[3]
Define a priority requirement.
\end{exercise}

\begin{exercise}[3]
Explain injury in a priority construction.
\end{exercise}

\begin{exercise}[3]
Explain finite injury.
\end{exercise}

\begin{exercise}[3]
Explain why finite injury does not mean finitely many total injuries across
an infinite construction.
\end{exercise}

\begin{exercise}[2]
Define the arithmetic hierarchy.
\end{exercise}

\begin{exercise}[2]
Define

\[
\Sigma^0_1.
\]
\end{exercise}

\begin{exercise}[2]
Define

\[
\Pi^0_1.
\]
\end{exercise}

\begin{exercise}[2]
Define

\[
\Delta^0_1.
\]
\end{exercise}

\begin{exercise}[3]
Explain why

\[
\Sigma^0_1
=
\text{c.e.}
\]
\end{exercise}

\begin{exercise}[3]
Explain why

\[
\Pi^0_1
=
\text{co-c.e.}
\]
\end{exercise}

\begin{exercise}[3]
Explain why

\[
\Delta^0_1
=
\text{computable}.
\]
\end{exercise}

\begin{exercise}[3]
Classify a formula of form

\[
\exists y\forall z\,R(x,y,z)
\]

when \(R\) is computable.
\end{exercise}

\begin{exercise}[3]
Classify a formula of form

\[
\forall y\exists z\,R(x,y,z)
\]

when \(R\) is computable.
\end{exercise}

\begin{exercise}[3]
Explain why bounded quantifiers do not add arithmetic-hierarchy levels.
\end{exercise}

\begin{exercise}[3]
State Post's theorem conceptually.
\end{exercise}

\begin{exercise}[3]
Explain the relationship between

\[
\Delta^0_2
\]

and computation relative to

\[
\mathbf{0}'.
\]
\end{exercise}

\begin{exercise}[3]
Explain what it means for a set to be

\[
\Sigma^0_n\text{-complete}.
\]
\end{exercise}

\begin{exercise}[3]
Explain why the arithmetic hierarchy is not a runtime hierarchy.
\end{exercise}

\begin{exercise}[2]
Define a computable approximation.
\end{exercise}

\begin{exercise}[3]
Construct the standard increasing approximation to a c.e.\ set.
\end{exercise}

\begin{exercise}[3]
Explain limit computability.
\end{exercise}

\begin{exercise}[3]
State the limit lemma conceptually.
\end{exercise}

\begin{exercise}[3]
Explain why a limit-computable approximation may change its answer before
stabilizing.
\end{exercise}

\begin{exercise}[3]
Explain why Turing degrees form a partial rather than linear order.
\end{exercise}

\begin{exercise}[3]
Define the join

\[
A\oplus B.
\]
\end{exercise}

\begin{exercise}[3]
Show that

\[
A\leq_TA\oplus B
\]

and

\[
B\leq_TA\oplus B.
\]
\end{exercise}

\begin{exercise}[3]
Explain enumeration reducibility conceptually.
\end{exercise}

\begin{exercise}[3]
Distinguish uniform and nonuniform computability statements.
\end{exercise}

\begin{exercise}[3]
Explain why

\[
\forall x\exists y\,P(x,y)
\]

does not automatically imply the existence of a computable choice
function.
\end{exercise}

\begin{exercise}[3]
Explain how recursion theory connects arithmetic definitions with
computational procedures.
\end{exercise}

\begin{exercise}[3]
Explain why theorem sets of effective formal theories are naturally c.e.
\end{exercise}

\begin{exercise}[4]
Develop primitive recursive coding of finite sequences.
\end{exercise}

\begin{exercise}[4]
Prove closure of primitive recursive functions under bounded minimization.
\end{exercise}

\begin{exercise}[4]
Study a proof that the Ackermann function is not primitive recursive.
\end{exercise}

\begin{exercise}[4]
Study equivalence of partial recursive and Turing-computable functions.
\end{exercise}

\begin{exercise}[4]
Study Kleene's normal form theorem.
\end{exercise}

\begin{exercise}[4]
Prove the \(s\)-\(m\)-\(n\) theorem.
\end{exercise}

\begin{exercise}[4]
Study acceptable numberings of partial computable functions.
\end{exercise}

\begin{exercise}[4]
Prove Kleene's recursion theorem.
\end{exercise}

\begin{exercise}[4]
Construct a self-reproducing program using the recursion theorem.
\end{exercise}

\begin{exercise}[4]
Study the recursion theorem with parameters.
\end{exercise}

\begin{exercise}[4]
Study creative and productive sets.
\end{exercise}

\begin{exercise}[4]
Prove a standard c.e.\ completeness theorem for the halting set.
\end{exercise}

\begin{exercise}[4]
Construct computably inseparable c.e.\ sets.
\end{exercise}

\begin{exercise}[4]
Study the construction of a simple set.
\end{exercise}

\begin{exercise}[4]
Study immune and hyperimmune sets.
\end{exercise}

\begin{exercise}[4]
Study Friedberg--Muchnik priority constructions conceptually.
\end{exercise}

\begin{exercise}[4]
Construct an intermediate c.e.\ degree using a finite-injury strategy at a
conceptual level.
\end{exercise}

\begin{exercise}[4]
Study incomparable Turing degrees.
\end{exercise}

\begin{exercise}[4]
Study low and high c.e.\ degrees.
\end{exercise}

\begin{exercise}[4]
Study jump inversion conceptually.
\end{exercise}

\begin{exercise}[4]
Prove basic closure properties of the arithmetic hierarchy.
\end{exercise}

\begin{exercise}[4]
Study Post's theorem in detail.
\end{exercise}

\begin{exercise}[4]
Study completeness of standard sets at each finite arithmetic-hierarchy
level.
\end{exercise}

\begin{exercise}[4]
Study the Shoenfield Limit Lemma.
\end{exercise}

\begin{exercise}[4]
Study enumeration degrees.
\end{exercise}

%%%%%%%%%%%%%%%%%%%%%%%%%%%%%%%%%%%%%%%%%%%%%%%%%%%%%%%%%%%%
\subsection{Conceptual Summary}
\label{subsec:recursion-theory-summary}
%%%%%%%%%%%%%%%%%%%%%%%%%%%%%%%%%%%%%%%%%%%%%%%%%%%%%%%%%%%%

Recursion theory develops the internal structure of computability.

Primitive recursive functions are generated from

\[
\boxed{
\text{zero}
+
\text{successor}
+
\text{projection}
}
\]

using

\[
\boxed{
\text{composition}
}
\]

and

\[
\boxed{
\text{primitive recursion}.
}
\]

They are all total, but they do not include every total computable function.

Adding unbounded minimization yields the partial recursive functions, which
coincide with partial Turing-computable functions.

Computable functions can be indexed:

\[
\boxed{
\varphi_0,
\varphi_1,
\varphi_2,\ldots.
}
\]

The \(s\)-\(m\)-\(n\) theorem formalizes effective parameter substitution,
while Kleene's recursion theorem gives program fixed points:

\[
\boxed{
\varphi_e
=
\varphi_{f(e)}.
}
\]

Computably enumerable sets can be written

\[
\boxed{
W_e
=
\operatorname{dom}(\varphi_e).
}
\]

Reducibility compares computational information, and Turing equivalence
produces degrees:

\[
\boxed{
\deg_T(A).
}
\]

The jump operation creates strictly stronger information:

\[
\boxed{
A<_TA'.
}
\]

The arithmetic hierarchy organizes definability into

\[
\boxed{
\Sigma^0_n,
\qquad
\Pi^0_n,
\qquad
\Delta^0_n.
}
\]

Post's theorem connects these logical levels with iterated jumps.

Priority methods reveal that c.e.\ degrees contain rich intermediate and
incomparable structure.

Recursion theory therefore links:

\[
\boxed{
\text{effective functions}
+
\text{program indices}
+
\text{oracle information}
+
\text{logical definability}
+
\text{degrees of unsolvability}.
}
\]

%%%%%%%%%%%%%%%%%%%%%%%%%%%%%%%%%%%%%%%%%%%%%%%%%%%%%%%%%%%%
\subsection{Section Connections}
\label{subsec:recursion-theory-connections}
%%%%%%%%%%%%%%%%%%%%%%%%%%%%%%%%%%%%%%%%%%%%%%%%%%%%%%%%%%%%

\chapterconnection{
Recursion Theory deepens Section~\ref{sec:computability-theory} by replacing
the basic computable-versus-uncomputable distinction with a detailed
structural theory of effective functions, c.e.\ sets, indices,
parameterization, fixed points, reducibilities, Turing degrees, jumps,
priority constructions, and the arithmetic hierarchy. It also connects
directly with Proof Theory because formal proofs and computation histories
can be coded arithmetically, and with Model Theory through definability and
effective structures. The next section, Type Theory, approaches foundations
from a different direction: mathematical expressions are assigned types,
proofs become structured terms, and the Curry--Howard correspondence turns
logical deduction into computation.
}

%%%%%%%%%%%%%%%%%%%%%%%%%%%%%%%%%%%%%%%%%%%%%%%%%%%%%%%%%%%%
% SECTION SOURCE TRACKING
%%%%%%%%%%%%%%%%%%%%%%%%%%%%%%%%%%%%%%%%%%%%%%%%%%%%%%%%%%%%

%------------------------------------------------------------
% 12.6 Recursion Theory
%------------------------------------------------------------
%
% Source basis:
%   - Standard classical recursion theory.
%   - Standard recursive-function theory.
%   - Standard c.e.-set, reducibility, degree, jump, priority, and
%     arithmetic-hierarchy theory.
%
% Used for:
%   - recursion theory versus computability theory;
%   - initial recursive functions;
%   - composition;
%   - primitive recursion;
%   - primitive recursive arithmetic;
%   - primitive recursive predicates;
%   - bounded search and bounded quantification;
%   - Ackermann function;
%   - limitations of primitive recursion;
%   - minimization;
%   - partial recursive functions;
%   - total recursive functions;
%   - effective program numberings;
%   - indices;
%   - universal partial computable functions;
%   - s-m-n theorem;
%   - parameterization;
%   - Kleene normal form;
%   - computation predicates;
%   - Kleene recursion theorem;
%   - program fixed points;
%   - self-reference;
%   - computably enumerable sets;
%   - W_e notation;
%   - stage approximations;
%   - co-c.e.\ sets;
%   - halting set;
%   - many-one completeness;
%   - index sets;
%   - productive sets;
%   - creative sets;
%   - computable inseparability;
%   - immune sets;
%   - simple sets;
%   - many-one reducibility;
%   - one-one reducibility;
%   - truth-table reducibility;
%   - Turing reducibility;
%   - Turing equivalence;
%   - Turing degrees;
%   - computable and halting degrees;
%   - Turing jump;
%   - iterated jumps;
%   - relativization;
%   - Post's problem;
%   - intermediate c.e.\ degrees;
%   - priority method;
%   - requirements and injury;
%   - finite-injury arguments;
%   - arithmetic hierarchy;
%   - Sigma, Pi, and Delta levels;
%   - bounded quantifiers;
%   - Post's theorem;
%   - hierarchy completeness;
%   - computable approximations;
%   - limit computability;
%   - limit lemma;
%   - degree joins;
%   - enumeration degrees preview;
%   - uniformity and effective choice;
%   - applications and exercises.
%
% Notes:
%   - Type Theory is developed next in Section 12.7.
%   - Computability Theory appears in Section 12.5.
%   - Constructive Mathematics follows in Section 12.8.
%   - Philosophy and Foundations follows in Section 12.9.
%
%%%%%%%%%%%%%%%%%%%%%%%%%%%%%%%%%%%%%%%%%%%%%%%%%%%%%%%%%%%%